% Pass geometry options before the class loads them:
%\PassOptionsToPackage{letterpaper, margin=1in}{geometry}

\documentclass[
  english,
  class=report,
  accentcolor=2c,
  %custommargins=geometry,  % Let geometry manage margins
  marginpar=false,
  parskip=half-,
  fontsize=11pt,
  ruledheaders=false,
  pdfa=false,
  IMRAD=false,
  thesis={type=master}
]{tudapub}

% -------------------------------------------------------------------
% Let TUDa load geometry once with our options:
%\usepackage{geometry}

% Math support (must load before anything uses \text, \binom, etc.):
\usepackage{amsmath}
\usepackage{amssymb}

% Language / microtype:
\usepackage[english, main=english]{babel}
\usepackage[autostyle]{csquotes}
\usepackage{microtype}
\usepackage{booktabs}
\usepackage{scrlayer-scrpage}       % If not already loaded
\usepackage{tikz}
\usetikzlibrary{positioning,arrows.meta,shapes.geometric,fit}
\usepackage{pgfplots}
\pgfplotsset{compat=1.18}
\usepgfplotslibrary{groupplots}
\usepackage{algorithm}
\usepackage{algpseudocode}
\usepackage{listings}
\lstset{
  basicstyle=\ttfamily\small,
  columns=fullflexible,
  keepspaces=true,
  breaklines=true,
  frame=single,
  numbers=left,
  numberstyle=\tiny,
  xleftmargin=1.8em,
  language=Python,
}
\usepackage{url}
\KOMAoptions{
  headsepline=false,  % or false if you want to remove the line entirely
  footsepline=false, % same for the footer
  headwidth=text,    % match the rule width to the text width
  footwidth=text
}
\pagestyle{scrheadings}

% -------------------------------------------------------------------
% 1) Times-like fonts (common in conference style)
%\usepackage{newtxtext,newtxmath}   % tudapub already sets up its own fonts — this clashes
%\usepackage[T1]{fontenc}    % ensure proper accent/symbol encoding
%\usepackage[utf8]{inputenc} % for UTF-8 input (not always needed nowadays)
% 2) Tighter heading layout with `titlesec`:

\setkomafont{chapter}{\Large\bfseries}
\setkomafont{section}{\large\bfseries}
\setkomafont{subsection}{\bfseries}
\setkomafont{subsubsection}{\bfseries}

\RedeclareSectionCommand[
  beforeskip=2ex plus .2ex,
  afterskip=1ex plus .2ex
]{chapter}
\renewcommand*\chapterformat{\thechapter.}

\RedeclareSectionCommand[
  beforeskip=1.5ex plus .2ex,
  afterskip=1ex plus .2ex
]{section}

\RedeclareSectionCommand[
  beforeskip=1.3ex plus .2ex,
  afterskip=0.8ex plus .2ex
]{subsection}

\RedeclareSectionCommand[
  beforeskip=1.2ex plus .2ex,
  afterskip=0.7ex plus .2ex
]{subsubsection}

% 3) Prevent vertical stretching:
\raggedbottom

\newenvironment{centeredabstract}[1][Abstract]{%
  \clearpage
  \thispagestyle{empty}
  \begin{center}
    {\Large \bfseries #1}\\[1em]
    \begin{minipage}{0.8\textwidth}
      \normalsize
}{%
    \end{minipage}
  \end{center}
  \clearpage
}

% -------------------------------------------------------------------
% Bibliography
\usepackage{biblatex}
\addbibresource{references.bib}

% -------------------------------------------------------------------
\begin{document}

\title{Graph-Conditioned Causal Data Augmentation: Reusable Local Dependencies for Synthetic Data Generation}
\author{Ramya Ravichandran}
\studentID{2729732}
\reviewer{Prof. Dr. Heinz Koeppl \and M. Sc. Philipp Fröhlich}

\department{etit}
\group{Self-Organizing Systems Lab}

\submissiondate{01.12.2026}

\maketitle
\affidavit

\begin{centeredabstract}
Causal discovery algorithms are usually evaluated and trained on a handful of small,
hand-curated datasets, simply because realistic data with a known ground-truth
graph is scarce. This thesis builds a tool that manufactures such data on
demand. Rather than sampling from a single fixed parametric structural causal
model, as most existing generators do, the approach learns a library of local
dependencies -- one bivariate copula per variable pair, and one small vine
copula per two-parent relationship -- directly from real datasets, and treats
this library as a set of reusable building blocks. Given any user-specified
directed acyclic graph, with a node indegree bounded by two in the core
setting and up to four in an extension, the sampler composes the corresponding
building blocks in topological order to synthesize a full dataset. The method is
validated on held-out data in four steps. First, the strongest fitted
pairwise dependencies on the Sachs protein-signaling data are predominantly
known causal edges (all nine strongest pairs on the observational data;
AUROC 0.75 against the consensus network). Second, resampling a graph with
the learned models reproduces its non-uniform per-edge dependence, where a
generic Gaussian baseline collapses every edge to the same strength.
Third, the learned building blocks transfer to graphs that never existed
in the source data: one-parent edges reproduce the dependence of the
copula they were built from to within $\Delta\tau \approx 0.007$ on average, a dependence
learned on one dataset can be transplanted onto an edge between variables
of another, and the generated data cannot be gamed by sorting variables by
their variance. Fourth, on the Asia network a structure-learning algorithm (PC) recovers
the true graph from the generated data significantly better than from a
fixed-correlation baseline or a fitted Gaussian copula, with no significant
difference from real data (which bounds, but does not rule out, a
difference), but a score-based algorithm (GES) does markedly worse on the
generated data than on real data, and on the larger Insurance and Alarm
networks both algorithms fall far short of real data, so this usefulness
does not yet extend to large discrete networks. In a prototype that
replaces the copula conditionals with random-forest conditionals fitted on
the same training data, both algorithms come within about three
structural errors of real data on Insurance and Alarm ($97$--$120\%$ of
the gap between the baseline and real data closed), at the price that these conditionals are
tied to one graph and dataset rather than being reusable building blocks. The
pipeline extends to discrete and mixed variables; continuizing categorical
columns by jittering dilutes Kendall's $\tau$ to its $\tau_a$ value, which
severely weakens dependence for skewed binary variables, and fitting the
copulas with the native discrete likelihood instead reduces the error of
reproduced edge dependence by up to two orders of magnitude on mixed data.
On discrete data, however, the learned copula generators do not reliably preserve
the conditional independences of the supplied graph: for chains and forks
with a discrete middle variable, the implied independence is rejected in
25--56\% of tests, against 5--7\% for real data, refitted conditional
probability tables and random-forest conditionals.
The result is a tested, reproducible generator of benchmark data for
user-specified graphs whose local dependencies are learned from real
data rather than assumed.
\end{centeredabstract}

\begin{centeredabstract}[Kurzfassung]
\begin{otherlanguage}{german}
Kausale Entdeckungsalgorithmen werden üblicherweise anhand einer Handvoll
kleiner, von Hand kuratierter Datensätze evaluiert und trainiert, einfach
weil realistische Daten mit bekanntem Ground-Truth-Graphen rar sind. Diese
Arbeit entwickelt ein Werkzeug, das solche Daten bei Bedarf erzeugt.
Anstatt aus einem einzigen fest vorgegebenen parametrischen strukturellen
Kausalmodell zu ziehen, wie es die meisten bestehenden Generatoren tun,
lernt der vorgestellte Ansatz eine Bibliothek lokaler Abhängigkeiten --
eine bivariate Copula pro Variablenpaar sowie eine kleine Vine-Copula pro
Zwei-Eltern-Beziehung -- direkt aus realen Datensätzen und behandelt diese
Bibliothek als eine Menge wiederverwendbarer Bausteine. Für einen
beliebigen, vom Nutzer vorgegebenen gerichteten azyklischen Graphen (DAG),
dessen Eingangsgrad der Knoten im Kernumfang auf zwei und in einer Erweiterung auf
bis zu vier begrenzt ist, setzt der Sampler die entsprechenden Bausteine in
topologischer Reihenfolge zusammen, um einen vollständigen Datensatz zu
synthetisieren. Die Methode wird auf zur\"uckgehaltenen Daten in vier Schritten
validiert. Erstens sind die st\"arksten angepassten paarweisen
Abh\"angigkeiten der Sachs-Daten \"uberwiegend bekannte kausale Kanten (alle
neun st\"arksten Paare der Beobachtungsdaten; AUROC 0{,}75 gegen\"uber dem
Konsensnetzwerk). Zweitens reproduziert das erneute Abtasten eines Graphen
mit den gelernten Modellen dessen ungleichm\"a{\ss}iges Abh\"angigkeitsmuster,
w\"ahrend eine generische Gau{\ss}-Baseline jede Kante auf dieselbe St\"arke
reduziert. Drittens lassen sich die gelernten Bausteine auf Graphen
\"ubertragen, die in den Ausgangsdaten nicht existieren: Kanten mit einem
Elternknoten reproduzieren die Abh\"angigkeit ihrer Copula im Mittel bis auf
$\Delta\tau \approx 0{,}007$, eine auf einem Datensatz gelernte Abh\"angigkeit
kann auf eine Kante zwischen Variablen eines anderen Datensatzes
\"ubertragen werden, und die erzeugten Daten lassen sich nicht durch
Sortieren nach der Varianz \"uberlisten. Viertens rekonstruiert der
PC-Algorithmus das Asia-Netzwerk aus den erzeugten Daten signifikant besser
als aus einer Baseline mit fester Korrelation oder einer angepassten
Gau{\ss}-Copula, ohne signifikanten Unterschied zu echten Daten (was einen
Unterschied eingrenzt, aber nicht ausschlie{\ss}t); ein scorebasierter
Algorithmus (GES) schneidet auf den erzeugten Daten jedoch deutlich
schlechter ab als auf echten Daten, und bei den gr\"o{\ss}eren Netzwerken
Insurance und Alarm bleiben beide Algorithmen deutlich hinter echten Daten
zur\"uck. Ersetzt man in einem Prototyp die Copula-Bedingungsmodelle durch
auf denselben Trainingsdaten angepasste Random-Forest-Modelle, erreichen
beide Algorithmen auf Insurance und Alarm bis auf etwa drei strukturelle
Fehler das Niveau echter Daten
($97$--$120\,\%$ der L\"ucke zwischen Baseline und echten Daten
geschlossen); diese Modelle sind jedoch an einen Graphen und Datensatz
gebunden und keine wiederverwendbaren Bausteine. Die Pipeline
wird auf diskrete und gemischte Variablen erweitert; das Kontinuisieren
kategorialer Spalten durch Jittering verd\"unnt Kendalls $\tau$ auf seinen
$\tau_a$-Wert, w\"ahrend die Anpassung mit der nativen diskreten Likelihood
den Fehler der reproduzierten Kantenabh\"angigkeit bei gemischten Daten um
bis zu zwei Gr\"o{\ss}enordnungen senkt. Bei diskreten Daten erhalten die
gelernten Copula-Generatoren die bedingten Unabh\"angigkeiten des vorgegebenen Graphen
jedoch nicht zuverl\"assig: F\"ur Ketten und Gabeln mit diskreter mittlerer
Variable wird die implizierte Unabh\"angigkeit in 25--56\,\% der Tests
verworfen, gegen\"uber 5--7\,\% f\"ur echte Daten, neu gesch\"atzte bedingte
Wahrscheinlichkeitstabellen und Random-Forest-Modelle. Das Ergebnis ist ein getesteter,
reproduzierbarer Generator von Benchmarkdaten f\"ur vom Nutzer vorgegebene
Graphen, dessen lokale Abh\"angigkeiten aus echten Daten gelernt statt
angenommen werden.
\end{otherlanguage}
\end{centeredabstract}

\KOMAoptions{DIV=15, BCOR=12mm}

\tableofcontents
\listoffigures
\listoftables
\listofalgorithms

\chapter*{List of Acronyms}
\addcontentsline{toc}{chapter}{List of Acronyms}
\begin{description}
  \item[AIC] Akaike Information Criterion
  \item[AUROC] Area Under the Receiver Operating Characteristic curve
  \item[BDeu] Bayesian Dirichlet equivalent uniform (score)
  \item[BIC] Bayesian Information Criterion
  \item[CDF] Cumulative Distribution Function
  \item[CI] Confidence Interval
  \item[CPDAG] Completed Partially Directed Acyclic Graph
  \item[CPT] Conditional Probability Table
  \item[DAG] Directed Acyclic Graph
  \item[gcda] Graph-Conditioned Causal Data Augmentation (this thesis's software package)
  \item[GES] Greedy Equivalence Search (score-based causal discovery)
  \item[KS] Kolmogorov--Smirnov (distance or test)
  \item[MEC] Markov Equivalence Class
  \item[MMD] Maximum Mean Discrepancy
  \item[NMI] Normalized Mutual Information
  \item[PC (algorithm)] Peter--Clark algorithm (constraint-based causal discovery)
  \item[SCM] Structural Causal Model
  \item[SHD] Structural Hamming Distance
  \item[TVD] Total Variation Distance
\end{description}

\chapter*{Notation}
\addcontentsline{toc}{chapter}{Notation}
\begin{description}
  \item[$x, y, \dots$] a raw variable / dataset column, on its original scale
  \item[$u, v \in (0,1)$] a pseudo-observation (rank-transformed variable), Equation~\ref{eq:pseudo-obs}
  \item[$C(\cdot,\cdot;\theta)$] a bivariate copula CDF with parameter(s) $\theta$
  \item[$c(\cdot,\cdot;\theta)$] the corresponding copula density
  \item[$\tau$] Kendall's rank correlation coefficient
  \item[$\rho$] linear (Pearson) or Gaussian-copula correlation parameter
  \item[$h_1(u_1,u_2;\theta)$] a copula's $h$-function, Equation~\ref{eq:hfunc}
  \item[$T(\cdot), T^{-1}(\cdot)$] the (inverse) Rosenblatt transform, Equation~\ref{eq:rosenblatt}
  \item[$G = (V,E)$] a directed acyclic graph with nodes $V$ and edges $E$
  \item[$\text{pa}(v)$] the parent set of node $v$ in a DAG
  \item[$n$] number of observations (rows) in a dataset or sample
  \item[$p$] number of variables (columns) in a dataset
  \item[$t_x$] the tie probability of a (discrete) variable $x$: the chance two independent draws of $x$ share a category
  \item[$\mathcal{L}$] a \texttt{DependencyLibrary} instance
\end{description}

\chapter{Introduction}
\label{chap:intro}

\section{Motivation}
Causal discovery methods -- algorithms that infer a directed acyclic graph
(DAG) of cause-and-effect relationships from purely observational data --
are usually benchmarked on a small, recycled set of datasets. The
protein-signaling network of \textcite{sachs2005} and the medical
Bayesian networks distributed with graphical-model toolkits such as
\emph{Asia} \parencite{lauritzen1988} and \emph{Insurance}
\parencite{binder1997} appear, in one form or another, in the majority of
papers on the topic. This is not because the field lacks appetite for more
data: it is because collecting a dataset with a \emph{known, verified} causal
graph attached to it is expensive, slow, or simply impossible for most
domains. Structure-learning methods that could otherwise benefit from more
training examples, more diverse graph topologies, or systematic stress-tests
against specific dependency strengths are left working with whatever handful
of graphs the literature happens to provide.

The usual workaround is to simulate data from a hand-specified structural
causal model (SCM): pick a DAG, attach a linear or additively-noisy equation
to every edge, and sample. This is cheap and flexible in principle, but it
throws away exactly the property that makes real benchmark datasets valuable
in the first place -- realistic, dataset-specific dependence structure. A
linear-Gaussian SCM cannot reproduce the fact that two proteins in a
signaling cascade are tied together by a strongly asymmetric, tail-dependent
relationship rather than a linear one, or that a binary risk factor and a
categorical outcome interact in a way no additive-noise model captures
cleanly.

This thesis asks whether the two approaches can be combined: keep the
flexibility of simulating for \emph{any} graph, but ground the local
dependencies in models that were actually fitted to real data.

\subsection{A Concrete Illustration}
Concretely, imagine a researcher developing a new constraint-independence
test for structure learning, who wants to check it specifically against
graphs with three-parent nodes -- a topology none of the standard
benchmarks exercises thoroughly: the Sachs consensus network has a single
three-parent node (Mek), Asia's largest indegree is 2, and Insurance reaches
indegree 3 at only five of its 27 nodes. Two options exist today. The first is to hand-write a
three-parent linear-Gaussian SCM: fast, but the resulting data has exactly
the generic dependence shape every other linear-Gaussian benchmark has,
so a good result says only that the new test handles \emph{that one
dependence shape} at that indegree, not that it handles three-parent nodes
in general. The second is to search for a real dataset with a known,
verified three-parent causal structure -- which, per the motivation above,
may simply not exist for the researcher's domain of interest. The tool
built in this thesis offers a third option: fit a dependency library from
whatever real datasets \emph{are} available (Sachs, UCI diabetes, Asia,
whatever domain-relevant data the researcher already has), then generate
data for the specific three-parent graph the researcher actually wants to
test against, with each edge's dependence grounded in a real fitted
relationship rather than a generic assumption. The sampler supports such
nodes: the baseline sampler hits its theoretical dependence up to six
parents, and learned conditional vines run without a single fallback up to
the four-parent nodes of the Alarm network (Section~\ref{sec:rq3}).

\section{Goal and Objective}
The goal of this thesis is a software tool that generates realistic
synthetic datasets for a user-specified graph structure by learning how
variables depend on each other from real data, rather than assuming a fixed
generic dependence model. Concretely, the tool
\begin{itemize}
  \item learns local dependencies -- primarily pairwise, and additionally
    conditional on up to two parents -- from real datasets, using copulas to
    separate each variable's marginal distribution from the way it co-varies
    with its neighbors, and
  \item reuses these learned dependencies as interchangeable building blocks
    to sample data for any DAG the user supplies, including graphs that
    combine building blocks learned from different source datasets and graph
    topologies that never occurred in any of them, with a node indegree
    bounded by a small constant (two in the core scope, extended to three or
    four).
\end{itemize}
Copulas are the natural tool for this separation: by Sklar's theorem
\parencite{sklar1959}, any multivariate distribution can be decomposed into
its marginals and a copula function that captures dependence alone, on a
common $(0,1)$ scale that is invariant to how each individual variable is
distributed. This is precisely the decomposition the tool needs, since it
must reuse a dependence pattern learned on one variable pair for a
structurally analogous pair with entirely different marginal statistics.

\section{Contributions}
Concretely, this thesis makes the following contributions:
\begin{enumerate}
  \item A reusable \emph{dependency library} abstraction that stores fitted
    pairwise and conditional (multi-parent) copula models as a validated,
    JSON-serializable object with provenance metadata, indexed by dataset
    and variable and decoupled from any particular graph; libraries learned
    on different datasets can be merged (Section~\ref{sec:pairwise-math} and
    Chapter~\ref{chap:method}).
  \item A graph sampler that consumes such a library and any user-specified
    DAG and produces a synthetic dataset by composing the corresponding
    building blocks in topological order: fitted vines where a parent set
    was observed, multi-parent models composed from pairwise copulas where
    it was not, copulas transplanted onto edges between variables of
    different datasets, and a counted, reported fallback for numerically
    unstable models (Section~\ref{sec:sampler}).
  \item An extension of the pipeline to discrete and mixed
    continuous/discrete variables, first by per-observation jittering with
    no change to the fitting or sampling code, and then by the native
    discrete copula likelihood, which removes most of the dependence loss
    caused by jittering (Sections~\ref{sec:discrete-method}
    and~\ref{sec:rq4}).
  \item A quantitative account of that dependence loss: jittering measures
    Kendall's $\tau_a$ rather than $\tau_b$, a known identity which here
    explains the per-edge degradation observed on Asia, Insurance and the
    mixed SCM (Section~\ref{sec:rq4}).
  \item A held-out, multi-seed evaluation with two structure-learning
    algorithms (PC and GES, 100 seeds per graph), showing that on the Asia
    network the learned generator's data supports both significantly better
    than a fixed-correlation baseline and a fitted Gaussian copula, and PC
    with no significant difference from real data, while on the larger
    Insurance and Alarm networks
    it falls clearly short of real data (Section~\ref{sec:closing-loop}).
  \item A demonstration, on a graph built from two datasets and on 200 novel
    graphs, that learned dependencies transfer to graph structures that
    never existed in any source dataset, reproducing the transplanted
    copulas and yielding benchmarks that are not variance-sortable
    (Section~\ref{sec:composition}).
  \item A check, on the chains, forks and colliders of four networks, of
    whether generated discrete and mixed data obey the conditional
    independences of the supplied graph. With a discrete middle variable the
    copula generators violate the chain and fork independences on the
    observed scale, while refitted CPTs and random-forest conditionals do
    not (Section~\ref{sec:rq4-ci}).
  \item A prototype of random-forest conditionals that work in the same
    pseudo-observation space and plug into the same sampler, which brings
    the structural error on Insurance and Alarm to within about three
    edges of that on real data and thereby supports the explanation that the remaining gap stems
    from non-monotone and interaction dependence that copula conditionals
    cannot represent (Sections~\ref{sec:forest-method}
    and~\ref{sec:rq6-forest}).
\end{enumerate}

\section{Scope}
The work is split into two phases. Phase 1 targets continuous variables
exclusively and establishes the full pipeline -- data loading, pairwise and
conditional dependency fitting, graph-conditioned sampling, and evaluation
against both statistical-realism and downstream-usability criteria. Phase 2
extends the same pipeline to binary, ordinal, and categorical variables,
without redesigning any of the Phase 1 machinery. Throughout, the core
target for graph complexity is two parents per node, with a controlled
extension to three and four parents to test how far the approach scales
before conditional vine fitting becomes unreliable. All experiments use a
combination of real datasets (Sachs, UCI diabetes, Asia, Insurance,
Alarm) and
synthetic ones with a known ground-truth DAG, so that recovered dependence
structure and downstream discovery performance can be checked against a
verifiable target rather than taken on faith.

\section{Thesis Structure}
Chapter~\ref{chap:background} introduces the statistical and causal-discovery
background the method builds on. Chapter~\ref{chap:method} describes the
pipeline itself and, importantly, the design decisions behind it --
several of which were not obvious from the tooling's documentation and had to
be reverse-engineered and validated before being trusted.
Chapter~\ref{chap:implementation} maps that design onto the actual codebase
and walks through two representative bugs at the code level.
Chapter~\ref{chap:setup} lays out the datasets, metrics, and experimental
protocol. Chapter~\ref{chap:results} answers six concrete research questions
with the corresponding experiments. Chapter~\ref{chap:discussion} reflects on
the bugs found along the way, the limitations that turned out to be
fundamental rather than fixable, and threats to validity.
Chapter~\ref{chap:conclusion} closes with a summary and directions for future
work.


\chapter{Background}
\label{chap:background}

\section{Causal Discovery and Structure Learning}
Causal discovery methods aim to recover a DAG over a set of variables from
observational data alone, typically by testing conditional independence
relations implied by the graph. The PC algorithm \parencite{spirtes2000} is
the canonical constraint-based method: it starts from a fully connected
undirected graph, removes an edge between two variables whenever some subset
of the other variables renders them conditionally independent, and then
orients the remaining edges using a small set of graphical rules. Its output
is generally a Markov equivalence class -- a completed partially directed
acyclic graph (CPDAG) -- rather than a single fully oriented DAG, since some
edges cannot be oriented from observational data alone regardless of how much
data is available.

Greedy Equivalence Search (GES) \parencite{chickering2002} is the canonical
score-based alternative. Instead of testing independences, it searches over
Markov equivalence classes: starting from the empty graph, it adds the
single edge that most improves a score of how well the graph fits the data,
until no addition helps, and then removes edges in the same greedy way. The
score rewards fit and penalizes complexity -- the Bayesian information
criterion (BIC) for continuous data, and for discrete data the Bayesian
Dirichlet equivalent uniform (BDeu) score \parencite{heckerman1995}, the
marginal likelihood of the data under a uniform prior on the conditional
probability tables. Like PC, GES returns a CPDAG. The two algorithms
therefore rely on the data in different ways: PC on the outcome of many
individual independence tests, GES on the overall fit of whole graphs.

The standard way to score how well a discovered graph matches the truth is
the Structural Hamming Distance (SHD): the number of edge insertions,
deletions, or reversals needed to turn one graph into the other, with zero
meaning an exact match. Formally, writing $E$ and $\hat{E}$ for the edge
sets (with direction) of the true and recovered graphs over the same
variable set,
\begin{equation}
\text{SHD}(E, \hat{E}) = |E \setminus \hat{E}| + |\hat{E} \setminus E| -
|\{(i,j) : (i,j) \in E,\, (j,i) \in \hat{E}\}|,
\label{eq:shd}
\end{equation}
i.e.\ the count of edges present in one graph but not the other, minus a
correction so that a reversed edge (present in both, but pointing opposite
ways) counts once rather than twice.

Figure~\ref{fig:mec-example} illustrates why SHD, rather than exact-match
accuracy, is the right metric here, and why even a perfect method can have
a nonzero SHD. The chain (a) and the common cause (b) encode the same
conditional-independence relations -- neither implies $A \perp C$
unconditionally, both imply $A \perp C \mid B$ -- and are therefore Markov
equivalent: no amount of observational data can tell a constraint-based
method such as PC which of the two generated the data, which is why PC
returns the completed partially directed acyclic graph (CPDAG) $A - B - C$
with both edges undirected for either of them. The collider (c) is
different: it implies $A \perp C$ marginally but \emph{not} $A \perp C \mid
B$, so it is distinguishable from (a) and (b) and PC can orient both of its
edges. If the true graph is the chain $A \to B \to C$ and a method returns
the equivalent common-cause structure $A \leftarrow B \to C$, only the
$A$--$B$ edge is reversed ($B \to C$ is identical), giving $\text{SHD} = 1$
by Equation~\ref{eq:shd} although both graphs make identical independence
claims. The SHD implementation used in this thesis (causal-learn) likewise
counts every undirected CPDAG edge whose true counterpart is directed as
one error -- a reminder that a low nonzero SHD, like the continuous-graph
results in Section~\ref{sec:closing-loop}, can reflect unavoidable
orientation ambiguity rather than a recovery failure.

\begin{figure}[ht]
\centering
\begin{tikzpicture}[
  node distance=1.0cm,
  var/.style={circle, draw, minimum size=7mm, inner sep=0pt, font=\small},
  lab/.style={rectangle, draw=none, inner sep=1pt, font=\footnotesize, anchor=north},
  arr/.style={-{Latex[length=2mm]}, thick}
]
\foreach \x/\i/\name in {0/1/{(a) chain}, 4.8/2/{(b) common cause}, 9.6/3/{(c) common effect (collider)}} {
  \node[var] (a\i) at (\x, 0) {$A$};
  \node[var, right=of a\i] (b\i) {$B$};
  \node[var, right=of b\i] (c\i) {$C$};
  \node[lab] at ([yshift=-0.5cm]b\i.south) {\name};
}
\draw[arr] (a1) -- (b1);
\draw[arr] (b1) -- (c1);
\draw[arr] (b2) -- (a2);
\draw[arr] (b2) -- (c2);
\draw[arr] (a3) -- (b3);
\draw[arr] (c3) -- (b3);
\end{tikzpicture}
\caption{Three DAGs with identical conditional-independence implications
except one: (a) and (b) both imply $A \perp C \mid B$ and are Markov
equivalent (indistinguishable from observational data); (c) does not imply
$A \perp C \mid B$ and is therefore distinguishable from the other two,
since conditioning on a collider \emph{induces} rather than removes
dependence between its parents.}
\label{fig:mec-example}
\end{figure}

SHD is the metric this thesis ultimately uses to answer whether synthetic
data is \emph{useful}, as opposed to merely \emph{realistic} -- a
distinction the experimental design in Chapter~\ref{chap:results} is built
around.

\section{Copulas}
A copula is a multivariate cumulative distribution function whose marginals
are all uniform on $(0,1)$. Sklar's theorem \parencite{sklar1959} states
that any joint distribution $F(x_1,\dots,x_p)$ can be written as
$F(x_1,\dots,x_p) = C(F_1(x_1),\dots,F_p(x_p))$ for some copula $C$, where
$F_i$ are the marginal distributions. This factorization is exactly what the
tool needs: it lets marginal shape (a protein's expression level, a
patient's BMI) be modeled and stored completely separately from dependence
structure, so that a dependence pattern learned from one dataset can, in
principle, be attached to a variable with entirely different marginal
statistics in another.

Bivariate copula families differ in the \emph{shape} of dependence they can
represent: the Gaussian copula is symmetric and has no tail dependence, the
Student-t copula adds symmetric tail dependence via a degrees-of-freedom
parameter, and Archimedean families such as Clayton, Gumbel, and Joe
\parencite{joe1997} capture asymmetric lower- or upper-tail dependence that
linear-Gaussian models cannot represent at all. Kendall's $\tau$
\parencite{kendall1938}, a rank-based dependence measure invariant to
monotonic transformations of the marginals, is the natural summary statistic
for comparing dependence strength across families and across the raw-versus-copula
divide, and is used throughout this thesis for exactly that reason.

For more than two variables, a \emph{vine copula} builds a full joint
dependence structure out of a cascade of bivariate copulas arranged in a tree
sequence, using conditional distributions computed via bivariate copulas'
$h$-functions. \textcite{nagler2016} give the general construction; this
thesis only needs the special case of a three-variable vine (one child, two
parents), used to sample a variable conditional on both of its parents at
once.

Figure~\ref{fig:vine-tree} shows the three-variable case exactly as it is
built in the code: a D-vine \parencite{bedford2002,aas2009,czado2019} whose
variables are listed in the order $C, P_1, P_2$, so that the first tree is
the \emph{path} $C - P_1 - P_2$. Tree $T_1$ contains the copulas of the
pairs $(C,P_1)$ and $(P_1,P_2)$; tree $T_2$ contains one copula for $C$
and $P_2$ \emph{conditional on} $P_1$, fitted not to the raw
pseudo-observations but to the $h$-function residuals $h(u_C \mid
u_{P_1})$ and $h(u_{P_2} \mid u_{P_1})$ from $T_1$. The joint density
factorizes as
\begin{equation}
c(u_C, u_{P_1}, u_{P_2}) = c_{C,P_1}(u_C, u_{P_1}) \cdot c_{P_1,P_2}(u_{P_1},
u_{P_2}) \cdot c_{C,P_2 \mid P_1}\bigl(h(u_C \mid u_{P_1}), h(u_{P_2}
\mid u_{P_1})\bigr),
\label{eq:vine-density}
\end{equation}
i.e.\ three bivariate copula densities multiplied together rather than one
trivariate density estimated directly -- the practical appeal of the vine
construction. Because $C$ is a leaf of the path, the conditional density
of the child given both parents is $c(u_C \mid u_{P_1}, u_{P_2}) =
c_{C,P_1}(u_C,u_{P_1}) \cdot c_{C,P_2\mid P_1}(h(u_C\mid u_{P_1}),
h(u_{P_2}\mid u_{P_1}))$: the parent--parent copula $c_{P_1,P_2}$ enters
the child's conditional distribution only through the conditioning
argument $h(u_{P_2}\mid u_{P_1})$. For $k$ parents the same construction
uses the path $C - P_1 - \dots - P_k$, again with the child as a leaf. A
D-vine and a C-vine coincide for three variables; for four or more they
differ, and only the D-vine path keeps the child a leaf in every tree,
which is what the sampling recipe of Section~\ref{sec:conditional} needs.

\begin{figure}[ht]
\centering
\begin{tikzpicture}[
  node/.style={circle, draw, minimum size=7mm, font=\small},
  edgelabel/.style={font=\scriptsize, fill=white, inner sep=1pt}
]
\node[node] (c) at (0,0) {$C$};
\node[node] (p1) at (2.4,0) {$P_1$};
\node[node] (p2) at (4.8,0) {$P_2$};
\draw (c) -- node[edgelabel] {$c_{C,P_1}$} (p1);
\draw (p1) -- node[edgelabel] {$c_{P_1,P_2}$} (p2);
\node[font=\small] at (2.4,0.7) {$T_1$};

\node[node, minimum size=9mm] (cp1) at (1.2,-2.2) {$C,P_1$};
\node[node, minimum size=9mm] (p1p2) at (3.6,-2.2) {$P_1,P_2$};
\draw (cp1) -- node[edgelabel] {$c_{C,P_2\mid P_1}$} (p1p2);
\node[font=\small] at (2.4,-1.5) {$T_2$};
\end{tikzpicture}
\caption{D-vine tree sequence used for a two-parent child $C$ (variable
order $C, P_1, P_2$). Tree $T_1$ is the path $C - P_1 - P_2$ with copulas
$c_{C,P_1}$ and $c_{P_1,P_2}$; tree $T_2$ joins its two edges with the
conditional copula $c_{C,P_2\mid P_1}$, fitted to the $h$-function
residuals of $T_1$. The child is a leaf, so it can be sampled conditional
on both parents by inverting the Rosenblatt transform
(Algorithm~\ref{alg:conditional}).}
\label{fig:vine-tree}
\end{figure}

\subsection{Copula Family Reference}
\label{sec:copula-glossary}
The candidate set offered to the AIC/BIC selection step in
Section~\ref{sec:pairwise-math} contains seven families: independence,
Gaussian, Student-$t$, Clayton, Gumbel, Frank and Joe. Clayton, Gumbel and
Joe additionally enter in their $90^\circ$, $180^\circ$ and $270^\circ$
rotations, so that negative dependence and the opposite tail can be
represented as well (12 of the 55 Sachs pairs select a $90^\circ$ or
$270^\circ$ rotation). The non-trivial families are defined here for
reference, since the results chapters name them without restating their
densities. All are written in terms of pseudo-observations $u,v \in
(0,1)$; the independence copula is $C(u,v) = uv$.

\textbf{Gaussian.} $C(u,v;\rho) = \Phi_2(\Phi^{-1}(u), \Phi^{-1}(v); \rho)$,
where $\Phi_2$ is the bivariate standard normal CDF with correlation $\rho
\in (-1,1)$ and $\Phi^{-1}$ the standard normal quantile function.
Symmetric, no tail dependence: extreme values in one variable carry no
extra information about extreme values in the other beyond what the bulk
of the distribution already implies.

\textbf{Student-t.} Same construction as Gaussian but with the
bivariate Student-$t$ CDF, adding a degrees-of-freedom parameter $\nu > 0$
that controls symmetric tail dependence in both directions -- lower $\nu$
means heavier joint tails. Reduces to the Gaussian copula as $\nu \to
\infty$.

\textbf{Clayton.} $C(u,v;\theta) = \bigl(u^{-\theta} + v^{-\theta} -
1\bigr)^{-1/\theta}$ for $\theta > 0$. Asymmetric: strong \emph{lower}-tail
dependence (joint small values are more likely than a Gaussian copula with
the same $\tau$ would predict) but no upper-tail dependence.

\textbf{Gumbel.} $C(u,v;\theta) = \exp\bigl(-\bigl[(-\log
u)^{\theta} + (-\log v)^{\theta}\bigr]^{1/\theta}\bigr)$ for $\theta \ge 1$.
The mirror image of Clayton: strong \emph{upper}-tail dependence, none in
the lower tail.

\textbf{Joe.} $C(u,v;\theta) = 1 - \bigl[(1-u)^{\theta} + (1-v)^{\theta}
- (1-u)^{\theta}(1-v)^{\theta}\bigr]^{1/\theta}$ for $\theta \ge 1$. Also
upper-tail dependent, with a different tail shape than Gumbel; it is Joe
that the pairwise fitting step in Section~\ref{sec:rq1} selects for
\emph{praf}--\emph{pmek} on the pooled Sachs data (Gumbel on the 853-row
observational subset).

\textbf{Frank.} $C(u,v;\theta) = -\frac{1}{\theta}\log\Bigl(1 +
\frac{(e^{-\theta u}-1)(e^{-\theta v}-1)}{e^{-\theta}-1}\Bigr)$ for $\theta
\ne 0$. Symmetric like Gaussian, but with \emph{no} tail dependence at all
in either direction, and can represent both positive and negative
dependence directly (Clayton, Gumbel and Joe cover negative dependence only
through their $90^\circ$/$270^\circ$ rotations). Frank is selected for all
three copulas of the Insurance conditional vine and for one of the three
Asia copulas in Table~\ref{tab:sweep}. Across \emph{all} pairwise fits,
however, independence, Clayton and Frank are the most frequent selections
on Insurance (149, 80 and 66 of 351 pairs), independence and Clayton on Asia
(14 and 6 of 28), and Frank is also the most frequent family on the
continuous UCI diabetes data (24 of 55 pairs); its prevalence therefore
cannot be attributed to jittering alone.

\section{Related Approaches to Synthetic Causal Data}
\label{sec:related-work}

\subsection{Structural-Causal-Model Simulators}
The dominant alternative to the approach in this thesis is to specify a
structural causal model directly: pick a DAG, and for every node write down a
function of its parents plus independent noise, most commonly a linear
combination with additive Gaussian noise. This is the generation method
behind most synthetic benchmarks used to evaluate continuous-optimization
structure-learning methods such as NOTEARS \parencite{zheng2018}, and it is
simple to implement and guarantees a known ground truth, which is why it
remains the default choice in benchmark suites for structure-learning
methods generally. Its weakness is exactly what this thesis addresses: every
edge in such a model has the same generic dependence shape regardless of
what the variables actually represent, so data generated this way is easy
for structure-learning methods to fit but does not stress-test them against
the kind of dependence patterns, such as tail asymmetry or bounded/discrete
support, that real data actually exhibits. A method that performs well on a
thousand linear-Gaussian SCMs has not necessarily been tested against
anything resembling the Sachs signaling network's own dependence structure,
which Section~\ref{sec:rq3} shows is predominantly non-Gaussian. Simulated
additive-noise benchmarks can moreover be ``gamed'': with generically
sampled edge weights the marginal variance tends to increase along the
causal order (high \emph{varsortability}), so that a trivial
sort-by-variance baseline competes with sophisticated methods
\parencite{reisach2021}. Section~\ref{sec:composition} checks this property
for the data generated in this thesis.

\subsection{Copula-Based Synthetic Data Generators}
Copula-based generation of synthetic tabular data is not new in itself: the
Synthetic Data Vault \parencite{patki2016} popularized fitting a copula (or a
Gaussian-mixture-copula hybrid) to an entire real table and resampling from
it, and the pure-Python \texttt{copulas} package this thesis considered and
rejected in favor of \texttt{pyvinecopulib} (Section~\ref{sec:pairwise-math}
below) grew out of that line of work. The key difference is what the fitted
dependence structure is \emph{for}. SDV-style tools fit a single joint model
to reproduce one table's own row-level statistics as faithfully as possible,
with the graph, if any, being incidental to the fitting procedure. This
thesis instead treats the pairwise and conditional copulas as detachable
building blocks, stored independently of any particular graph, so the same
learned dependence can be attached to a different edge, in a different DAG,
possibly built from a different dataset entirely -- the composability
demonstrated in Section~\ref{sec:composition}. Copula-based graphical
models with one copula per local family of a DAG have been proposed for
density estimation \parencite{elidan2010}; none of the generators surveyed
here exposes the recombination of fitted local dependencies across graphs
and datasets as a first-class, evaluated capability, which is the
methodological contribution this thesis aims at.

\subsection{Benchmark Suites and Toolkits for Causal Discovery}
On the discovery side, toolkits such as the Causal Discovery Toolbox
\parencite{kalainathan2020} and the \texttt{bnlearn} ecosystem
\parencite{scutari2010} -- the latter being the origin of the Asia and
Insurance networks used throughout this thesis -- standardize how discovery
algorithms are run and scored, but they consume benchmark datasets as given
rather than providing a way to manufacture new ones with controllable
dependence realism. This thesis's contribution sits upstream of that
tooling: \texttt{gcda}'s output is deliberately just another CSV that any of
these downstream toolkits, causal-learn included, can consume unmodified.

\subsection{Deep Generative Tabular Synthesizers}
A fourth line of work generates synthetic tabular data with deep generative
models rather than copulas: CTGAN \parencite{xu2019} adapts a conditional
GAN to mixed continuous/categorical tables, using a mode-specific
normalization scheme to handle multimodal continuous columns and a
training-by-sampling procedure to handle imbalanced categorical ones. These
methods can capture dependence structure copulas cannot easily express
(high-order, nonlinear interactions among many variables at once) and have
no explicit per-pair or per-edge structure to speak of -- which is
precisely their limitation for this thesis's purposes: a trained CTGAN
model reproduces one table's joint distribution as a monolithic whole, with
no natural notion of a reusable, edge-level building block, and no
mechanism for attaching a subset of its learned dependence to a different,
user-specified graph. They also typically need substantially more training
rows than copula fitting to train reliably, which cuts against small
benchmark datasets such as Sachs (a few thousand rows) or the synthetic
SCM used throughout this thesis.

\subsection{Realistic Benchmark Generators Fitted to Real Data}
\label{sec:realistic-generators}
The closest prior work fits a generator to real data and uses it to produce
benchmark datasets with a known ground truth. RealCause
\parencite{neal2020} fits generative models to real treatment--outcome data
so that causal-effect estimators can be benchmarked on realistic data with
known effects. CausalTime \parencite{cheng2024} fits neural models with
normalizing flows to real multivariate time series, derives a hypothesised
causal graph and generates time series from the fitted model.
causalAssembly \parencite{goebler2024} uses a production line whose causal
structure is known from the physical process, estimates each node's
conditional distribution with distributional random forests and samples
from the resulting joint distribution, which is Markov with respect to the
ground truth. All three reproduce the \emph{source system's own} graph (or
subgraphs of it); none of them is designed to transplant a fitted local
dependence onto a different graph or to combine local models fitted on
different datasets, which is the focus of this thesis.

\subsection{Comparison Summary}
Table~\ref{tab:related-comparison} summarizes where each line of work
stands on the three properties Section~\ref{sec:related-work} identifies as
jointly necessary: learning dependence from real data rather than assuming
a fixed form, exposing that dependence as building blocks reusable across
graphs, and validating against a downstream discovery task rather than
distributional realism alone.

\begin{table}[ht]
\caption{Comparison against related approaches to synthetic data for causal-discovery-adjacent use cases.}
\label{tab:related-comparison}
\centering
\small
\begin{tabular}{p{2.6cm}p{2.3cm}p{2.9cm}p{3.1cm}}
\toprule
\textbf{Approach} & \textbf{Learns from real data} & \textbf{Composable building blocks} & \textbf{Downstream-validated} \\
\midrule
Linear-Gaussian SCM \parencite{zheng2018} & No & Yes (trivially: fixed form) & Sometimes \\
SDV / copula-fit table \parencite{patki2016} & Yes & No & No \\
CTGAN / deep generative \parencite{xu2019} & Yes & No & Rarely \\
RealCause / CausalTime / causalAssembly \parencite{neal2020,cheng2024,goebler2024} & Yes & No (source graph only) & Yes (benchmarks) \\
\textbf{This thesis} & \textbf{Yes} & \textbf{Yes} & \textbf{Yes} (Section~\ref{sec:closing-loop}) \\
\bottomrule
\end{tabular}
\end{table}

\subsection{Positioning of This Thesis}
Put together, the gap this thesis targets is the combination none of the
above lines of work covers on its own: dependence learned from real
data (unlike SCM simulators), organized as reusable building blocks rather
than a single fitted table or trained model (unlike SDV-style generators
and deep generative synthesizers alike), and evaluated against the actual
downstream discovery task rather than only against realism metrics (unlike
most copula-fitting work, which stops at distributional comparison).
Chapter~\ref{chap:results} is structured specifically to demonstrate all
three properties rather than assume they follow from each other.


\chapter{Method}
\label{chap:method}

\section{Pipeline Overview}
The pipeline has two stages that are deliberately kept independent: a
\emph{fitting} stage that turns a real dataset into a library of reusable
dependency models, and a \emph{sampling} stage that consumes any such
library (or several, merged) together with a user-specified DAG and produces
a new dataset. Figure~\ref{fig:pipeline} sketches the data flow. Keeping the
two stages separate, with a serializable library object as the sole
interface between them, is what makes cross-dataset composition in
Section~\ref{sec:composition} possible: the sampler never needs to know
which dataset a given pairwise model came from.

\begin{figure}[ht]
\centering
\begin{tikzpicture}[
  box/.style={draw, rounded corners, minimum width=2.6cm, minimum height=1cm, align=center, font=\small},
  arr/.style={-{Latex[length=2mm]}, thick}
]
\node[box] (data) {Real dataset};
\node[box, right=1.4cm of data] (prep) {Preprocessing\\(pseudo-obs.)};
\node[box, right=1.4cm of prep] (fit) {Pairwise +\\conditional fitting};
\node[box, right=1.4cm of fit] (lib) {Dependency\\Library (JSON)};
\node[box, below=1.2cm of lib] (graph) {User-specified\\DAG};
\node[box, below=1.2cm of fit] (sampler) {GraphSampler};
\node[box, below=1.2cm of prep] (out) {Synthetic\\dataset};
\node[box, below=1.2cm of data] (eval) {Evaluation};

\draw[arr] (data) -- (prep);
\draw[arr] (prep) -- (fit);
\draw[arr] (fit) -- (lib);
\draw[arr] (lib) -- (sampler);
\draw[arr] (graph) -- (sampler);
\draw[arr] (sampler) -- (out);
\draw[arr] (out) -- (eval);
\draw[arr] (data) -- (eval);
\end{tikzpicture}
\caption{Pipeline overview. Fitting (top row) and sampling (bottom row)
communicate only through the serialized dependency library, so a graph can
be sampled using models learned from a different dataset than the one it
was originally observed in, or from several datasets merged together.}
\label{fig:pipeline}
\end{figure}

\section{Preprocessing: Pseudo-Observations}
Copula fitting requires each column to be transformed to a uniform
$(0,1)$ scale first. Formally, for a column of $n$ observations
$x_1,\dots,x_n$, the pseudo-observation for $x_i$ is
\begin{equation}
\hat{u}_i = \frac{\operatorname{rank}(x_i)}{n+1},
\label{eq:pseudo-obs}
\end{equation}
where $\operatorname{rank}(x_i) \in \{1,\dots,n\}$ is $x_i$'s position in
the sorted column. Fitting copulas to such rank-based pseudo-observations by maximum likelihood
is the standard semiparametric approach \parencite{genest1995}. Using
$(n+1)$ rather than $n$ as the denominator, the
standard plotting-position convention for empirical CDFs, guarantees
$\hat{u}_i \in (0,1)$ strictly, so no pseudo-observation ever lands exactly
on the boundary $0$ or $1$, where several copula densities (Clayton's and
Joe's among them) are undefined or infinite. The transform is wrapped in an
\texttt{EmpiricalMarginal} object that also stores the inverse mapping
$\hat{u} \mapsto x$ via linear interpolation between order statistics, so
that samples drawn on the copula scale can be mapped back to the original
variable's units after generation.

This preprocessing step produced the most consequential bug of the project.
An early version of the transform function accepted no explicit argument for
which values to transform and defaulted to transforming the array's own
\emph{sorted} order rather than the caller's original row order. Every
pairwise fit downstream was therefore silently computed on scrambled data.
The bug was not caught by code inspection; it surfaced only because a sanity
check reported Kendall's $\tau \approx 0.999$ for every variable pair,
regardless of which variables were involved -- a value no real dataset
should plausibly produce across the board. This is worth dwelling on
because it illustrates a general point about this kind of pipeline: a silent
data-alignment bug does not throw an exception, and generic unit tests that
check shapes and types would not have caught it. What caught it was a
domain-informed sanity check on the \emph{output distribution} of a
statistic that should vary. The fix was to require every caller to pass the
original values explicitly, removing the ambiguity, and a regression test
now checks that transforming a permuted copy of a column and then inverting
it recovers the original order.

\section{Pairwise Dependency Fitting}
\label{sec:pairwise-math}
For every unordered pair of variables in a dataset, a bivariate copula
density $c(\hat{u},\hat{v};\theta)$ is fit to the pair's pseudo-observations
by maximum likelihood, separately for each candidate family $\mathcal{F} \in
\{\text{independence}, \text{Gaussian}, \text{Student-}t, \text{Clayton},
\text{Gumbel}, \text{Frank}, \text{Joe}\}$, the last three families also in
their rotated versions (Section~\ref{sec:copula-glossary}):
\begin{equation}
\hat{\theta}_{\mathcal{F}} = \arg\max_{\theta} \sum_{i=1}^{n} \log
c_{\mathcal{F}}(\hat{u}_i, \hat{v}_i; \theta).
\label{eq:mle}
\end{equation}
The family itself is then selected by whichever of Akaike Information
Criterion (AIC) or Bayesian Information Criterion (BIC) is being used for
that experiment,
\begin{align}
\text{AIC}_{\mathcal{F}} &= 2k_{\mathcal{F}} - 2\log
L_{\mathcal{F}}(\hat{\theta}_{\mathcal{F}}), \\
\text{BIC}_{\mathcal{F}} &= k_{\mathcal{F}} \log n - 2\log
L_{\mathcal{F}}(\hat{\theta}_{\mathcal{F}}),
\label{eq:aicbic}
\end{align}
where $k_{\mathcal{F}}$ is the number of free parameters in family
$\mathcal{F}$ (zero for independence; one for Gaussian, Clayton, Gumbel,
Frank and Joe; two for Student-$t$) and $L_{\mathcal{F}}$ is the
likelihood at the fitted parameters. The library builder uses BIC by
default; Section~\ref{sec:rq3} compares it with AIC. The family minimizing
the chosen criterion is kept, along with its rotation, its parameters, its
Kendall's $\tau$, and
both $\text{AIC}_{\mathcal{F}}$ and $\text{BIC}_{\mathcal{F}}$, as
JSON-serializable metadata for that pair -- so a later re-analysis (such as
the AIC-versus-BIC comparison in Section~\ref{sec:rq3}) does not require
refitting anything. \texttt{pyvinecopulib} was chosen as the underlying
library over the pure-Python \texttt{copulas} package because it exposes
this exact AIC/BIC family selection natively -- needed regardless for the
conditional case below -- has a C++ backend for the larger sweeps, and
includes vine-copula support out of the box.

Before trusting the library for anything beyond a toy example, its core
fit/sample round trip was validated on data with a known ground truth:
Gaussian, Student-t, and Clayton copulas with fixed parameters were sampled,
refit, and checked to recover the correct family, parameters, and $\tau$;
samples drawn from the \emph{fitted} (not the true) copula were then checked
to match the same dependence statistics as the original data. All three
families passed, with one expected wrinkle: the Student-t copula's
degrees-of-freedom parameter converges far more slowly than its correlation
parameter and was still 7.6\% off ($\hat\nu = 4.30$ for $\nu = 4$) at the
$n = 3{,}000$ used by the validation script (\path{t3_validation.py}), which
is a known property of that parameter rather than a fitting defect, and is
handled by
giving it a looser relative tolerance in the validation test than the
absolute tolerance used for the correlation parameter.

\section{Conditional Sampling for Two-Parent Nodes}
\label{sec:conditional}
A node with two parents cannot be sampled from a single bivariate copula;
it needs a distribution for the child conditional on both parents at once.
This is implemented with a three-variable D-vine and the inverse Rosenblatt
transform. For a bivariate copula $C(u_1,u_2;\theta)$, the $h$-function
\begin{equation}
h_1(u_1,u_2;\theta) = \frac{\partial C(u_1,u_2;\theta)}{\partial u_1} =
P(U_2 \le u_2 \mid U_1 = u_1)
\label{eq:hfunc}
\end{equation}
gives the conditional distribution of one variable given the other, and its
inverse $h_1^{-1}(u_1,w;\theta)$ maps a fresh uniform draw $w$ to a sample
from that conditional distribution. For the path-shaped D-vine of
Figure~\ref{fig:vine-tree} (variable order $C, P_1, P_2$), the forward
Rosenblatt transform \parencite{rosenblatt1952} as computed by
\texttt{pyvinecopulib} maps a joint sample to independent uniforms,
processing the variables from the end of the path,
\begin{equation}
T(u_C, u_{P_1}, u_{P_2}) = \bigl(h(u_C \mid u_{P_1}, u_{P_2}),\;
h(u_{P_1} \mid u_{P_2}),\; u_{P_2}\bigr),
\label{eq:rosenblatt}
\end{equation}
with $h(u_C \mid u_{P_1}, u_{P_2}) = h_{C \mid P_2; P_1}\bigl(h(u_C \mid
u_{P_1}) \,\big|\, h(u_{P_2} \mid u_{P_1})\bigr)$ computed from the copulas
$c_{C,P_1}$, $c_{P_1,P_2}$ and $c_{C,P_2\mid P_1}$. The two parent
coordinates of $T$ do not depend on $u_C$, which is what makes the
following recipe valid. Sampling a leaf of a D-vine conditional on the
remaining variables in this way is standard, e.g.\ in D-vine quantile
regression \parencite{kraus2017}; what had to be established empirically
before it could be trusted was \texttt{pyvinecopulib}'s ordering
convention. The recipe for sampling $C \mid P_1{=}p_1, P_2{=}p_2$ is: forward-transform the triple $(\cdot, p_1, p_2)$, with a
placeholder in the child's slot, through Equation~\ref{eq:rosenblatt};
overwrite that transformed placeholder with a fresh draw $w \sim
\text{Unif}(0,1)$; and apply the inverse Rosenblatt transform $T^{-1}$ to
obtain a joint sample whose child coordinate is the desired conditional
draw. Listing the child \emph{first} in the vine's variable order is what
makes it the fully-conditioned variable in $T$'s processing order --
confirmed empirically by fixing one column at a time and observing which
transformed coordinates changed as a result. Algorithm~\ref{alg:conditional}
summarizes the procedure. Before using this in the pipeline, it was checked
against the closed-form conditional mean and variance of a known trivariate
Gaussian correlation matrix and matched to within Monte Carlo noise. The
same conditional-sampler component is reused, unmodified, for the three-parent
extension tested in Section~\ref{sec:sweep}, with the corresponding
closed-form Gaussian check committed as a regression test.

\begin{algorithm}
\caption{Conditional sampling of a two-parent child via a D-vine}
\label{alg:conditional}
\begin{algorithmic}[1]
\Require fitted D-vine with child $C$ listed first, parent values $p_1, p_2 \in (0,1)$
\Ensure a draw $u_C$ from $C \mid P_1{=}p_1, P_2{=}p_2$
\State $w \gets$ fresh draw from $\text{Unif}(0,1)$
\State $(\cdot, v_1, v_2) \gets T(w, p_1, p_2)$ \Comment{forward Rosenblatt, Eq.~\ref{eq:rosenblatt}; child slot is a placeholder}
\State $(u_C, p_1', p_2') \gets T^{-1}(w, v_1, v_2)$ \Comment{inverse Rosenblatt with parent coordinates unchanged}
\State \Return $u_C$
\end{algorithmic}
\end{algorithm}

The direction of the $h$-function inverse used for the simpler one-parent
case raised the same kind of ambiguity -- which of $h_1(u_1,u_2)$'s two
arguments is the conditioning value -- and was resolved the same way: a
large sample from a Clayton copula with a known parameter was binned near a
fixed conditioning value, and the empirical conditional quantiles of the
other variable were compared against the library's output before the
function was trusted for real data. This check is now part of the test
suite
(\path{test_hinv1_direction_matches_empirical_conditional_quantiles_of_clayton}).

\section{The Dependency Library}
The dependency library is the interface between fitting and sampling: a
JSON-serializable container holding every pairwise model, every fitted
conditional (two-parent) model, and every marginal transform, indexed by
dataset and variable name. Making this an explicit, serializable object
rather than passing fitted model objects around in memory was a deliberate
design choice: it is what allows two libraries fit on entirely different
datasets to be merged into one before sampling
(\texttt{DependencyLibrary.merge}, which refuses colliding dataset
namespaces), which is the mechanism behind the cross-dataset composition
experiment in Section~\ref{sec:composition}. Every library also stores
provenance metadata (package versions, git commit, arguments) and a schema
version, and is validated when loaded.

Reconstructing a stored copula from its serialized parameters surfaced a
second bug. For a pair where the independence family was selected -- no
dependence at all -- the fitted copula's parameter array has shape
$(0,0)$; converting it to a plain list and back with \texttt{numpy} collapses
that to shape $(0,)$, which the copula constructor rejects. This went
unnoticed through several demos because none of them \emph{sampled} a
one-parent edge whose stored copula was the independence family:
independence was selected for six UCI-diabetes pairs and one synthetic-SCM
pair, but none of these pairs was an edge of a sampled graph. It surfaced
when reconstructing the full Asia network in Phase 2, which does contain
such edges. It was fixed with an explicit reshape guard and a regression test
covering the independence-family case specifically -- a reminder that a
pipeline validated end-to-end on a handful of examples has, at best,
validated those examples, not the space of cases the data can produce.

\section{Graph Sampler}
\label{sec:sampler}
Given a DAG (a \texttt{networkx.DiGraph}, indegree bounded at two by default
and configurably higher) and a dependency library, \texttt{GraphSampler}
draws an $n \times p$ dataset one node at a time in topological order.
All nodes are generated on the copula scale and mapped to their own units at
the end, through the library's fitted marginal or a standard normal baseline
marginal if none is given; a root node (indegree zero) is therefore an
independent uniform draw. Indegree-one nodes are sampled via the
corresponding pairwise copula's $h$-function inverse, falling back to a
fixed-correlation Gaussian copula if no learned model exists for that pair;
if the node has a conditional model instead of a pair copula (a random
forest, Section~\ref{sec:forest-method}), that model is used. Indegree-two-or-more nodes
use the conditional vine sampler from Section~\ref{sec:conditional} when a
vine has been fitted for exactly that parent set (looked up by the
\emph{set} of parents, so the order in which a user lists the edges does
not matter). For a parent set that never occurred together in a source
dataset, the child's conditional density is instead composed from the
learned pairwise copulas of its incoming edges,
$c(u_C \mid u_{P_1},\dots,u_{P_k}) \propto \prod_j c_{C,P_j}(u_C, u_{P_j})$,
and sampled by numerical inversion of its CDF on a fixed grid; this is
exact when the parents are conditionally independent given the child in
copula space and otherwise an approximation (it is the local density of a
copula Bayesian network, \cite{elidan2010}). If neither is available, a
Gaussian vine with correlation $\rho$ between the child and each parent is
used as the baseline; since that correlation matrix is positive definite
only if $k\rho^2 < 1$, $\rho$ is capped at $\sqrt{0.9/k}$ (unchanged for
$\rho = 0.5$ and $k \le 3$). Algorithm~\ref{alg:sampler} lays out the full
loop, including where each fallback triggers.

\begin{algorithm}
\caption{\texttt{GraphSampler}: sampling a dataset for an arbitrary DAG (as implemented in \texttt{GraphSampler.sample})}
\label{alg:sampler}
\begin{algorithmic}[1]
\Require DAG $G$, for every node $v$ a marginal $F_v$ (standard normal if none) and an optional edge model from the library $\mathcal{L}$; sample size $n$; clipping bound $\varepsilon$; baseline correlation $\rho$
\Ensure synthetic dataset $D$ with $n$ rows
\For{node $v$ in topological order of $G$}
  \State $P \gets$ parents of $v$ in $G$, \quad $k \gets |P|$
  \State \textbf{try}
  \If{$k = 0$}
    \State $u_v \gets n$ draws from $\text{Unif}(\varepsilon, 1-\varepsilon)$
  \ElsIf{$k = 1$}
    \If{$v$ has a conditional model but no pair copula} \Comment{e.g.\ a random forest}
      \State $u_v \gets$ draw from that model given $u_P$
    \ElsIf{$v$ has a pair copula $c$ (learned or bound)}
      \State $u_v \gets h_c^{-1}(w \mid u_P)$, $w \sim \text{Unif}(\varepsilon,1-\varepsilon)$, in $c$'s column orientation (Eq.~\ref{eq:hfunc})
    \Else
      \State $u_v \gets h^{-1}(w \mid u_P)$ of a Gaussian copula with correlation $\rho$ \Comment{baseline}
    \EndIf
  \Else
    \If{$v$ has a conditional model $M$ for the parent set $P$} \Comment{vine, composed, or forest}
      \State $u_v \gets$ draw from $M$ given $u_P$ in $M$'s parent order (Algorithm~\ref{alg:conditional} for a vine)
    \Else
      \State baseline: Gaussian vine, correlation $\min(\rho, \sqrt{0.9/k})$ to each parent; draw via Algorithm~\ref{alg:conditional}
    \EndIf
  \EndIf
  \State \textbf{on exception:} $u_v \gets n$ draws from $\text{Unif}(\varepsilon,1-\varepsilon)$; fallback count of $v \gets n$
  \State replace non-finite entries of $u_v$ by $\text{Unif}(\varepsilon,1-\varepsilon)$ draws and count them
  \State $u_v \gets \text{clip}(u_v, \varepsilon, 1-\varepsilon)$
  \If{$k > 0$}
    \If{\texttt{calibrate\_margins}} \State $u_v \gets \text{rank}(u_v)/(n+1)$
    \ElsIf{$\text{KS}(u_v, \text{Unif}(0,1)) > \max(0.05,\, 1.95/\sqrt{n})$} \State emit a marginal-drift warning
    \EndIf
  \EndIf
\EndFor
\State \Return $D$ with column $x_v = F_v^{-1}(u_v)$ for every node $v$
\end{algorithmic}
\end{algorithm}

Every node's sampling step is wrapped in a try/except block with an explicit
non-finite check, and falls back immediately to independent uniform draws
rather than retrying: for the whole column if the step raises an exception,
and for the affected rows only if some draws are not finite. This design choice -- fail fast and move
on, rather than retry -- was made deliberately: a single unstable fit (for
instance a near-degenerate conditional vine on a highly skewed discrete
variable, the exact case surfaced in Phase 2) should not be allowed to abort
generation of an otherwise-usable $n \times p$ dataset, and a silent retry
loop would obscure exactly how often instability occurs. Instead, every
fallback draw is counted per node (\texttt{GraphSampler.fallback\_counts\_})
and announced with a dedicated warning class, which is what makes the
stability comparisons in Section~\ref{sec:discrete-results} possible: the
number of fallback draws is a measurable quantity, not a hidden
implementation detail.

A second diagnostic concerns compatibility. A multi-parent model is fitted
under the parent distribution of its source data; if the parents generated
in a new graph are distributed differently, the child's generated uniforms
are no longer uniform and mapping them through the stored marginal does not
reproduce that marginal. The sampler therefore measures, for every node,
the Kolmogorov--Smirnov (KS) distance of its generated uniforms from $U(0,1)$
and warns when it exceeds both 0.05 and the 0.1\% critical value. With the
option \texttt{calibrate\_margins}, each node's uniforms are replaced by
their normalized ranks before use: ranks, and hence the generated copula,
are unchanged, but every node's marginal is reproduced exactly.

\section{Worked Example: Sampling a Three-Node Graph}
\label{sec:worked-example}
To make Algorithms~\ref{alg:conditional} and~\ref{alg:sampler} concrete
rather than abstract, this section traces one full sample through a
three-node graph $A \to C$, $B \to C$ by hand, with illustrative (not
data-derived) numbers chosen to be easy to follow rather than realistic.

\textbf{Setup.} Suppose the dependency library holds: a Gaussian marginal
$\mathcal{N}(0,1)$ for $A$, a Gaussian marginal $\mathcal{N}(5,2)$ for $B$,
a Clayton copula with $\theta=2$ (Kendall's $\tau \approx 0.5$) for the pair
$(A,C)$, a Gaussian copula with $\rho=0.3$ ($\tau \approx 0.19$) for the
pair $(B,C)$, and a conditional vine for $C \mid A, B$ with the path
structure of Figure~\ref{fig:vine-tree} ($C - A - B$). For simplicity
assume its first tree uses a Clayton copula with $\theta=2$ for $(C,A)$ and
independence for $(A,B)$ (the two roots are independent), and its second
tree a Gaussian copula with $\rho=0.3$ for $(C,B \mid A)$. Note that the
vine's copulas are selected jointly when the vine is fitted
(\texttt{Vinecop.select}); they are not copied from the pairwise library,
which is used only for one-parent nodes.

\textbf{Step 1 -- root nodes ($A$, $B$).} Both have indegree zero, so
Algorithm~\ref{alg:sampler} draws them directly from their marginals. Say
this draw produces $a = 0.84$ (a slightly-above-average draw from
$\mathcal{N}(0,1)$) and $b = 6.1$ (slightly-above-average from
$\mathcal{N}(5,2)$). Their pseudo-observations, via each marginal's CDF,
are $u_A = \Phi(0.84) \approx 0.80$ and $u_B \approx \Phi((6.1-5)/2)
\approx 0.71$.

\textbf{Step 2 -- the two-parent node ($C$).} $C$ has indegree two, so
Algorithm~\ref{alg:sampler} delegates to Algorithm~\ref{alg:conditional}.
A fresh uniform draw, say $w = 0.62$, takes the place of $C$'s first
Rosenblatt coordinate $h(u_C \mid u_A, u_B)$ (Equation~\ref{eq:rosenblatt}).
Inverting the transform for this coordinate with the parents held fixed
first inverts the second tree, $v = h^{-1}_{C\mid B; A}\bigl(w \mid
h(u_B\mid u_A)\bigr)$, which with independent $A,B$ gives $h(u_B \mid u_A) =
u_B = 0.71$, and then the first tree, $u_C = h^{-1}_{C\mid A}(v \mid u_A)$.
Because the Clayton copula linking $C$ and $A$ is the stronger dependence
($\tau \approx 0.5$ versus $\tau \approx 0.19$ for the second-tree Gaussian
copula), $u_C$ is pulled further toward $u_A = 0.80$ than toward $u_B$ --
an asymmetric pull determined by each fitted copula rather than a fixed
blend, which is exactly what Section~\ref{sec:rq2} shows a
Gaussian-everywhere baseline cannot reproduce, since one fixed correlation
on every edge weights both parents identically.

\textbf{Step 3 -- inverse marginal transform.} Whatever pseudo-observation
$u_C \in (0,1)$ results is finally mapped back to $C$'s own units via its
\texttt{EmpiricalMarginal} inverse (or a standard normal quantile if $C$
has no library marginal), completing one row of the synthetic dataset.
Repeating Steps 1--3 $n$ times, independently, produces the full $n \times
3$ table.

\textbf{Where fallbacks bite.} If the $T_2$ copula's fitted parameters
happened to be numerically degenerate for this particular draw of $w$ --
the scenario Section~\ref{sec:discrete-results} shows is common for highly
skewed discrete children -- Algorithm~\ref{alg:sampler}'s non-finite check
would catch the resulting \texttt{NaN} and substitute an independent
$\text{Unif}(0,1)$ draw for $u_C$ instead, incrementing $C$'s fallback
counter but allowing generation of the remaining $n-1$ rows (and the rest
of the graph, in a larger example) to proceed unaffected.

\section{Extension to Discrete and Mixed Data}
\label{sec:discrete-method}
Extending the pipeline to binary, ordinal, and categorical variables was
approached by continuizing discrete columns rather than building a
parallel discrete-specific pipeline. Every discrete column is mapped to
$(0,1)$ by \emph{jittering}: each individual observation of a given category
receives an independent uniform draw on that category's slice of the
empirical CDF, rather than a single shared rank for the whole category. This
matters because a shared rank would make every observation of, say,
``smoker = yes'' identical on the copula scale, silently destroying any
information a copula fit could use and making the transform non-invertible
per observation; the per-observation jitter keeps it invertible and keeps
ties from collapsing the pseudo-observations onto a handful of points. Once
continuized, a discrete column is indistinguishable, as far as any copula
fitting or sampling code is concerned, from a genuinely continuous one --
which is exactly the point: none of Sections~\ref{sec:conditional}
through~\ref{sec:sampler} needed to change to support discrete data.

For variables whose categories have a natural order -- Insurance's
socio-economic status (\emph{Prole} $<$ \emph{Middle} $<$ \emph{UpperMiddle}
$<$ \emph{Wealthy}) is one of several such cases -- the category order used
for the jittering slices must reflect that true order rather than
alphabetical order, since fitting a copula against the wrong order corrupts
concordance for any pair involving that variable. This was missed in an
initial pass and caught during an explicit audit against the original task
plan (Section~\ref{sec:audit}); the fix threads each dataset's true
category order through to the jittering transform. \texttt{Asia}'s eight
binary variables \parencite{lauritzen1988} and \texttt{Insurance}'s
twenty-seven categorical variables with a maximum indegree of three
\parencite{binder1997} are both sampled from their bundled conditional
probability tables via \texttt{pgmpy} \parencite{ankan2015}, which ships the
same networks distributed by the reference source as local files and needs
no network access.

The jittering approach is a documented, deliberately simple choice; it is
a continuous-convolution technique that is well founded for many estimation
tasks \parencite{nagler2018}, but for concordance it has a known cost
\parencite{denuit2005,genest2007}: pairs that are tied in the original
data are ordered at random after jittering and contribute zero expected
concordance, so Kendall's $\tau$ of the jittered data equals, in
expectation, Kendall's $\tau_a$ of the raw data rather than $\tau_b$
(Section~\ref{sec:rq4}). As an alternative, \texttt{pyvinecopulib} can fit
copulas to discrete data natively: the likelihood of an observation is the
probability of the rectangle $[F(x^-), F(x)] \times [F(y^-), F(y)]$ under
the copula, which needs $F(x)$ and $F(x^-)$ but no jittering. This mode is
implemented as \texttt{discrete\_mode="native"} for both pairwise copulas
and conditional vines (the fitted latent copula is then sampled exactly as
in the continuous case) and compared with jittering in
Section~\ref{sec:rq4}. Since copulas of discrete data are not unique
\parencite{genest2007}, the native fit estimates the latent copula of a
continuous variable discretized at the observed category boundaries.

\section{Random-Forest Conditionals}
\label{sec:forest-method}
A copula conditional, whether an inverse $h$-function or a vine, represents
dependence that is monotone along each variable's order -- for a
categorical variable, along its fixed category order -- and approximates
interactions between parents through the simplifying assumption.
Conditional probability tables of discrete networks are bound by neither.
On the training draws of Asia, Insurance and Alarm, 0 of 6,
3 of 34 and 7 of 40 strong edges (Cram\'er's $V > 0.2$) have a
Spearman correlation of the category indices below half their Cram\'er's
$V$, and several tables are logical functions of the parents, such as Asia's
\emph{either} = \emph{tub} OR \emph{lung}.

As a nonparametric alternative for a node's conditional given its parents,
a random forest \parencite{breiman2001} is fitted in the same
pseudo-observation space as the rest of the pipeline. It therefore plugs
into the graph sampler unchanged: the parents enter with their
$U$-values, and the child's $U$-value is mapped back through its stored
marginal. For a discrete child, a classification forest predicts the
category probabilities; a category is drawn from them and turned into a
$U$-value uniform on that category's slice of the CDF. For a continuous
child, a regression forest is used as a conditional sampler in the manner of
quantile regression forests \parencite{meinshausen2006}: a tree is chosen
at random, and the sampler returns the $U$-value of a random training
observation from the leaf into which the parent values fall. Leaves hold at
least 10 (discrete) or 20 (continuous) training observations, so the
sampler does not replay individual rows. Unlike the copula library, a
forest is fitted for one node, one parent set and one dataset. It is
therefore not a reusable building block for new graphs or across datasets,
and is evaluated only on the known graphs of RQ6
(Section~\ref{sec:rq6-forest}).


\chapter{Implementation}
\label{chap:implementation}

Chapter~\ref{chap:method} described the pipeline's design; this chapter
describes how that design maps onto the actual codebase, and walks through
two of the bugs mentioned only in passing above in enough detail to show
what a fix in this kind of pipeline actually looks like at the code level.

\section{Module Structure}
Figure~\ref{fig:module-diagram} shows how the core packages under
\texttt{gcda/} depend on each other. Two further modules,
\texttt{gcda.experiments} (the
held-out, multi-seed experiments of Chapter~\ref{chap:results}) and
\texttt{gcda.provenance} (run metadata attached to every output), sit on top
of all of them and are omitted from the figure. The dependency direction is
deliberately one-way: \texttt{generation} (the sampler) depends on
\texttt{dependency} (the library) but not vice versa, and neither depends
on \texttt{evaluation}, so that evaluation code can be changed freely
without any risk of it being (accidentally) load-bearing for generation
itself -- a real risk in smaller research codebases where evaluation
scripts and core logic tend to accrete into the same files.

\begin{figure}[ht]
\centering
\begin{tikzpicture}[
  box/.style={draw, rounded corners, minimum width=3.1cm, minimum height=1cm, align=center, font=\small},
  arr/.style={-{Latex[length=2mm]}, thick}
]
\node[box] (data) {\texttt{gcda.data}\\loaders + preprocessing};
\node[box, right=1.6cm of data] (dep) {\texttt{gcda.dependency}\\pairwise + conditional fitting,\\DependencyLibrary};
\node[box, right=1.6cm of dep] (gen) {\texttt{gcda.generation}\\GraphSampler};
\node[box, below=1.3cm of dep] (eval) {\texttt{gcda.evaluation}\\realism, conditioning,\\discrete metrics, causal\_discovery};
\node[box, below=1.3cm of gen] (cli) {\texttt{gcda.cli}\\build-library / generate /\\evaluate / sweep};

\draw[arr] (data) -- (dep);
\draw[arr] (dep) -- (gen);
\draw[arr] (data) -- (eval);
\draw[arr] (gen) -- (eval);
\draw[arr] (dep) -- (cli);
\draw[arr] (gen) -- (cli);
\draw[arr] (eval) -- (cli);
\end{tikzpicture}
\caption{Module dependency graph. Arrows point from a module to the module that depends on it. \texttt{evaluation} depends on both \texttt{data} (for ground truth) and \texttt{generation} (for the data being evaluated), but nothing depends on \texttt{evaluation} in turn, keeping it safely swappable.}
\label{fig:module-diagram}
\end{figure}

\section{Case Study: The Pseudo-Observation Scrambling Bug}
\label{sec:bug-case-study-1}
Section~\ref{sec:pairwise-math}'s preprocessing step needs every column
mapped to pseudo-observations in the \emph{original} row order, since a
pairwise fit compares $\hat{u}_i$ and $\hat{v}_i$ at the same row $i$. The
first implementation of \path{EmpiricalMarginal.transform} looked
approximately like Listing~\ref{lst:bug1-before}: it computed ranks
internally but returned them in whatever order \texttt{numpy.argsort}
produced internally, rather than the caller's original order, whenever no
explicit values were passed in.

\begin{lstlisting}[caption={Before (simplified reconstruction for illustration; the buggy version is documented in the lab notebook but was not preserved in the repository history): silently returns pseudo-observations in sorted, not original, row order when called without an explicit argument.}, label={lst:bug1-before}]
class EmpiricalMarginal:
    def __init__(self, values):
        self.sorted_values = np.sort(values)
        self.n = len(values)

    def transform(self, values=None):
        # BUG: falls back to self.sorted_values, discarding
        # the caller's original row order entirely.
        target = values if values is not None else self.sorted_values
        ranks = np.searchsorted(self.sorted_values, target) + 1
        return ranks / (self.n + 1)
\end{lstlisting}

Every call site in the original pairwise-fitting loop happened to call
\texttt{transform()} with no argument, on the (false) assumption that the
object already "remembered" the right values from construction. The fix,
shown in Listing~\ref{lst:bug1-after}, removes the fallback entirely and
makes the argument mandatory, so the same mistake cannot recur silently at
a new call site in the future.

\begin{lstlisting}[caption={After (verbatim from \texttt{gcda/data/preprocessing.py}, docstrings shortened): \texttt{transform} has no default argument, so every caller must pass the values, in the order it wants them transformed.}, label={lst:bug1-after}]
@dataclass
class EmpiricalMarginal:
    sorted_values_: np.ndarray = field(default=None, repr=False)
    plotting_positions_: np.ndarray = field(default=None, repr=False)

    def fit(self, x):
        x = np.asarray(x, dtype=float)
        n = len(x)
        self.sorted_values_ = np.sort(x)
        self.plotting_positions_ = np.arange(1, n + 1) / (n + 1)
        return self

    def transform(self, x):
        x = np.asarray(x, dtype=float)
        return np.interp(x, self.sorted_values_, self.plotting_positions_)
\end{lstlisting}

The accompanying regression test, added alongside the fix, checks the
property that actually matters rather than any specific numeric output:
transforming a permuted copy of a column and inverting the result must
recover the original values in the original order, for an arbitrary
permutation.

\section{Case Study: The Independence-Family Reshape Bug}
Section~\ref{sec:pairwise-math} noted that reconstructing a stored copula
whose selected family is \texttt{indep} crashes on a shape mismatch.
Listing~\ref{lst:bug2-before} shows the deserialization code that
triggered it: the independence family's parameter array is
\texttt{numpy}-shaped $(0,0)$ before serialization, but round-tripping it
through \texttt{.tolist()} and back through \texttt{np.array(...)}
collapses the empty array to shape $(0,)$, which the underlying
\texttt{pv.Bicop} constructor rejects as malformed.

\begin{lstlisting}[caption={Before (simplified from \texttt{\_bicop\_from\_pairwise\_entry}): round-tripping an empty (0,0)-shaped parameter array through JSON collapses it to shape (0,), which the copula constructor rejects.}, label={lst:bug2-before}]
def _bicop_from_pairwise_entry(entry):
    family = getattr(pv.BicopFamily, entry["family"])
    params = np.array(entry["parameters"], dtype=float)  # BUG: (0,0) -> (0,)
    return pv.Bicop(family=family, rotation=entry["rotation"], parameters=params)
\end{lstlisting}

\begin{lstlisting}[caption={After (verbatim from \texttt{gcda/generation/graph\_sampler.py}): an explicit reshape guard restores the (0,0) shape the constructor expects whenever the family carries no parameters; the family name is also validated against pyvinecopulib's families before use.}, label={lst:bug2-after}]
def _bicop_from_pairwise_entry(entry: dict) -> pv.Bicop:
    name = entry["family"]
    if not isinstance(name, str) or not hasattr(pv.BicopFamily, name) or name.startswith("_"):
        raise ValueError(f"unknown copula family {name!r} in library entry")
    family = getattr(pv.BicopFamily, name)
    params = np.array(entry["parameters"], dtype=float)
    if params.size == 0:
        # e.g. the "indep" family has no parameters; .tolist() loses the (0, 0)
        # shape pv.Bicop's constructor requires, collapsing it to a bare [].
        params = params.reshape(0, 0)
    return pv.Bicop(family=family, rotation=int(entry["rotation"]), parameters=params)
\end{lstlisting}

This bug is a good illustration of a broader point from
Chapter~\ref{chap:discussion}: it was latent in the codebase for the
entire continuous-only Phase 1, because no one-parent edge sampled in
Phase 1 used an independence copula (independence was selected for six
UCI-diabetes pairs and one synthetic-SCM pair, none of them a sampled
edge), so \texttt{\_bicop\_from\_pairwise\_entry} was exercised only on
families with nonempty parameter arrays. It surfaced the moment Phase 2's Asia network,
which does contain independence-selected pairs, was reconstructed from a
saved library -- an example of a code path that passed every existing test
simply because nothing had ever asked it to handle that input.

\section{Testing Philosophy}
The regression tests added throughout Chapter~\ref{chap:method} (including
the discrete extension, Section~\ref{sec:discrete-method}) follow one
guiding rule, made explicit here because Chapter~\ref{chap:discussion}
argues it generalizes beyond this
specific project: whenever a bug is found because a summary statistic
looked implausible -- too uniform, too flat, too clean -- the regression
test asserts on that same statistic's plausible \emph{range} or
\emph{direction of change}, not on a specific fitted number. A test that
merely reproduces today's fitted parameters to five decimal places would
have passed throughout the entire period the pseudo-observation bug was
present, since the bug's symptom was a suspicious pattern across many
pairs, not a wrong-looking number in any single one.

\chapter{Experimental Setup}
\label{chap:setup}

\section{Datasets}
Table~\ref{tab:datasets} summarizes the datasets used. They were chosen to
cover three axes that matter for the research questions in
Chapter~\ref{chap:results}: real versus synthetic (so that recovered
structure can be checked against either domain knowledge or an exactly
known ground truth), continuous versus discrete, and small versus larger
graphs.

\begin{table}[ht]
\caption{Datasets used across the two phases of the pipeline.}
\label{tab:datasets}
\centering
\begin{tabular}{lp{3.1cm}p{6.7cm}}
\toprule
\textbf{Dataset} & \textbf{Type} & \textbf{Role} \\
\midrule
Sachs \parencite{sachs2005} & real, continuous & 11 phosphoproteins and phospholipids measured in 7,466 single cells pooled over nine experimental conditions (one general stimulation and eight interventions); the 853-row general-stimulation condition is used as the observational subset. Checks whether fitted pairwise models recover the consensus network. \\
UCI diabetes \parencite{efron2004} & real, continuous & 442 patients, 10 baseline variables and a disease-progression target; checks recovered dependence against its documented strongest predictors. \\
Synthetic SCM & synthetic, continuous & Small DAG with a known, deliberately varied edge-strength pattern; the only dataset where ``did sampling reproduce the true per-edge dependence'' can be checked exactly. \\
Asia \parencite{lauritzen1988} & real (CPT-sampled), discrete & 8 binary variables; the small sanity-check network for the discrete extension. \\
Insurance \parencite{binder1997} & real (CPT-sampled), discrete & 27 categorical variables, indegree up to 3; the larger discrete benchmark, including genuinely ordinal variables. \\
Alarm \parencite{beinlich1989} & real (CPT-sampled), discrete & 37 discrete variables, indegree up to 4; the task plan's alternative larger benchmark, used for the held-out discovery experiment and the indegree-4 check. \\
Discretized Sachs & synthetic (from real), discrete & Quantile-binned continuous data; a stress test with realistic (non-CPT) discrete dependence. \\
Mixed synthetic SCM & synthetic, mixed & Small DAG with a continuous$\to$discrete edge and a two-parent node with one continuous and one discrete parent; a shape neither Asia nor Insurance exercises on its own. \\
\bottomrule
\end{tabular}
\end{table}

\section{Evaluation Metrics}
Three complementary families of metrics are used, because none of them
alone answers the question this thesis cares about.

\emph{Statistical realism.} \texttt{compare\_marginals} reports the
Kolmogorov--Smirnov statistic and a decile-quantile $L_1$ distance between
true and simulated marginals; \path{compare_pairwise_dependence} reports
Kendall's $\tau$ and Spearman's $\rho$ per edge. A single multivariate
summary, the energy distance \parencite{szekely2013} and a Gaussian-kernel
maximum mean discrepancy \parencite{gretton2012}, is computed on
rank-transformed (U-space) data specifically so that it reflects
dependence-structure realism rather than being dominated by whichever
variable happens to have the largest raw scale. The held-out experiments
use the unbiased (U-statistic) estimators, which are centred at zero when
both samples come from the same distribution, together with a permutation
two-sample test (200 permutations, both samples subsampled to 500 rows).

\emph{Graph-conditioning checks.} \path{edge_removal_effect} verifies
that deleting an edge from the sampling graph collapses the corresponding
empirical dependence towards zero, and \path{dependence_strength_sweep}
verifies that increasing an edge's copula parameter increases the resulting
Kendall's $\tau$ monotonically. These checks target a different failure mode
than realism metrics: a sampler could produce individually realistic
marginals and pairwise dependence values while completely ignoring the
requested graph structure, and these two checks are what would catch that.

\emph{Discrete-specific metrics.} Category frequency comparison and joint
contingency distance serve as sanity checks; bias-corrected Cram\'er's V
\parencite{bergsma2013} and
normalized mutual information serve as the discrete analogue of Kendall's
$\tau$ for dependence strength; and \path{measure_generation_stability}
records wall-clock time, the number of fallback events (warnings of the
dedicated fallback class only) and the number of fallback draws from the
sampler's per-node counters, treating stability itself as a metric rather
than an implementation detail.

\emph{Downstream usability.} Structural Hamming Distance between the graph
recovered by the PC algorithm \parencite{spirtes2000} or GES
\parencite{chickering2002}, both run via
\texttt{causal-learn} \parencite{zheng2024}, and the true DAG, computed for
data generated by each method and for the real data itself, complemented by
the number of edges the algorithm finds and the precision, recall and $F_1$ score of
its skeleton (adjacencies) against the true skeleton. This is the
metric that most directly answers the thesis's stated goal, and is treated
as the headline result in Chapter~\ref{chap:results} rather than a secondary
check.

\section{Hyperparameter Sweep Design}
\label{sec:sweep}
A single sweep script varies three factors for the two-parent conditional
model: the copula family set offered to AIC/BIC selection (Gaussian-only
versus the full parametric set), the information criterion used for
selection (AIC versus BIC), and the dataset (all three continuous and both
discrete datasets). A separate sweep varies only the sampler's indegree
cap (2, 3, or 4 parents) to test how far conditional-vine sampling scales
before it becomes unreliable; it now covers caps 2 to 6, each repeated with
ten seeds, and uses the validity-capped baseline correlation of
Section~\ref{sec:sampler}. A third sweep varies the clipping epsilon of the
conditional sampler ($10^{-10}$ to $10^{-2}$), as requested by the task
plan's ``clipping/fallback settings''. What is swept is defined entirely in
a versioned YAML configuration file rather than hardcoded into the sweep
script, so that a specific sweep run can be reproduced from the config file
alone; every run additionally logs the clipping epsilon used during
sampling and the number of fallback draws it triggered, and its full
configuration is logged to Weights \& Biases alongside its metrics, in
offline mode by default so that reproducing a sweep locally requires no
external credentials.

\section{Protocol for the Held-Out Experiments}
\label{sec:reproduce-seeds}
The original evaluation fitted each library on one data draw and compared
generated data with true data from four seeds, one of which (seed 1 of the
synthetic SCM) was the library's own training data. All experiments added
in this version use a held-out protocol instead: the library is fitted once
on a training draw (seed 10{,}000), and every evaluation seed $s = 0, 1,
\dots$ draws independent true data and an independent generated sample.
The number of evaluation seeds is 100 for RQ2 and RQ6 (synthetic SCM,
Asia, Insurance, Alarm), 50 for RQ4 (Asia, Insurance, Alarm, mixed SCM) and
10 for the learned-vine indegree check; the novel-graph study uses 200
random DAGs. These runs were executed on the Lichtenberg cluster
(\path{slurm/}, Python 3.13); every seed they share with earlier runs on a
4-core workstation reproduced the workstation results exactly. Per seed,
PC ($\alpha = 0.05$; Fisher-$z$ for continuous, $\chi^2$ for discrete
data) and, as a second, score-based algorithm, GES (BIC score for
continuous, BDeu for discrete data) are run on the true data and on data from
every generator: the fixed-$\rho$ Gaussian baseline with real marginals,
a fitted Gaussian copula (the same pipeline restricted to Gaussian and
independence families), the learned generator, for discrete networks the
learned generator with native discrete fitting, and refitted conditional
probability tables (CPT refit), the natural generator for a purely discrete
Bayesian network. Results are reported as means with 95\% bootstrap
confidence intervals, and generators are compared with paired Wilcoxon
signed-rank and paired $t$-tests over seeds. The library-fitting variance
(a single training draw per experiment) is not resampled.
A non-significant difference is not read as equivalence. No formal
equivalence test with a pre-specified margin was performed; where a
generator is compared with real data, the 95\% confidence interval of the
paired difference is reported instead, and its bounds state how large a
difference the results are still compatible with. Phrases such as ``no
significant difference from real data'' are meant in exactly this sense.

\chapter{Results}
\label{chap:results}

The results are organized around six research questions, each targeting a
different claim the thesis goal depends on: that dependence can be learned
faithfully (RQ1), that it can be recomposed for an arbitrary graph (RQ2),
that the choice of copula family set actually matters and is
dataset-dependent (RQ3), that the same machinery extends to discrete data
(RQ4), that the learned dependencies are genuinely reusable across datasets
(RQ5), and, finally, that the resulting data is useful for the task that
motivated building it at all (RQ6).

\section{RQ1: Does the fitted library recover known real-world structure?}
\label{sec:rq1}
RQ1 asks whether strong fitted pairwise dependencies coincide with known
causal edges. All 55 Sachs variable pairs were ranked by the absolute
Kendall's $\tau$ of their fitted pairwise copula and the ranking was scored
against the skeleton of the Sachs consensus network as distributed by
bnlearn/pgmpy (17 edges). Because the bundled Sachs file pools nine
experimental conditions (Chapter~\ref{chap:setup}), the analysis is run both
on all 7,466 rows and on the 853-row general-stimulation (observational)
condition (the first 853 rows of the bundled file; verified row by row
against an independent copy of the observational data, 853 of 853 rows
identical). The ranking is scored by the area under the receiver operating
characteristic curve (AUROC): the probability that a randomly chosen
consensus edge is ranked above a randomly chosen non-edge, 0.5 for a random
ranking and 1 for a perfect one. Table~\ref{tab:rq1} summarizes the result.

\begin{table}[ht]
\caption{RQ1: ranking the 55 Sachs pairs by $|\tau|$ of the fitted pairwise copula against the 17-edge consensus skeleton. Random ranking gives precision $17/55 = 0.31$ and AUROC 0.5; the $p$-value is the one-sided Mann--Whitney test that consensus pairs have larger $|\tau|$.}
\label{tab:rq1}
\centering
\begin{tabular}{lcccccc}
\toprule
\textbf{Rows} & \textbf{AUROC} & \textbf{$p$} & \textbf{P@5} & \textbf{P@10} & \textbf{P@17} & \textbf{Rank of praf--pmek} \\
\midrule
Pooled (7,466) & 0.641 & 0.050 & 1.00 & 0.80 & 0.53 & 1 \\
Observational (853) & 0.752 & $1.1\times10^{-5}$ & 1.00 & 0.90 & 0.71 & 3 \\
\bottomrule
\end{tabular}
\end{table}

On the observational condition the nine strongest pairs are all consensus
edges, and the ranking separates edges from non-edges clearly (AUROC 0.75).
On the pooled data the separation is weaker (AUROC 0.64, $p = 0.05$): the
interventions shift several variables jointly and create strong dependence
between pairs that are not adjacent in the consensus graph. The
\emph{praf}--\emph{pmek} pair is the strongest pair on the pooled data
(fitted Joe copula, $\tau = 0.666$) and the third strongest on the
observational data (Gumbel, $\tau = 0.483$), behind
\emph{p44/42}--\emph{pakts473} and \emph{PKC}--\emph{P38}, both of which are
also consensus edges. Two caveats apply to the value 0.666: it is the
$\tau$ implied by the fitted copula, whereas the empirical Kendall's
$\tau_b$ of the raw data is 0.616, and 91\% of the \emph{praf} values are
ties because the measurements are rounded (the fitted copula is the same
under maximum-rank, mid-rank and random tie handling). On UCI diabetes,
\emph{s5} ($\tau = 0.41$) and \emph{bmi} ($0.36$) are the two variables most
strongly dependent on the progression target, the two predictors that
enter first in the analysis of \textcite{efron2004}.

Figure~\ref{fig:sachs-cascade} shows the relevant part of the consensus
network. Pairwise dependence is not causation -- the copula fit has no
access to edge directions, and strong dependence also arises between
non-adjacent variables -- so RQ1 establishes that the library captures the
\emph{strongest} real dependencies, which are predominantly direct edges,
not that it identifies the graph.

\begin{figure}[ht]
\centering
\begin{tikzpicture}[
  node/.style={draw, rounded corners, minimum width=1.8cm, minimum height=0.8cm, font=\small},
  arr/.style={-{Latex[length=2mm]}, thick}
]
\node[node] (praf) {\emph{praf}};
\node[node, right=1.3cm of praf] (pmek) {\emph{pmek}};
\node[node, right=1.3cm of pmek] (p4442) {\emph{p44/42}};
\node[node, right=1.3cm of p4442] (pakts) {\emph{pakts473}};
\draw[arr] (praf) -- node[above, font=\scriptsize] {$\tau=0.666$} (pmek);
\draw[arr] (pmek) -- (p4442);
\draw[arr] (p4442) -- node[above, font=\scriptsize] {$\tau=0.637^{*}$} (pakts);
\end{tikzpicture}
\caption{The raf/mek/erk cascade subgraph of the Sachs consensus network
(bnlearn/pgmpy). Edge labels: Kendall's $\tau$ of the fitted pairwise
copula on the pooled data; $^{*}$ on the observational subset, where
\emph{p44/42}--\emph{pakts473} is the strongest of all 55 pairs.}
\label{fig:sachs-cascade}
\end{figure}

\section{RQ2: Does the sampler reproduce learned dependence for an arbitrary graph?}
\label{sec:rq2}
Sampling the synthetic SCM's own DAG using the library learned from its own
data reproduces both the true marginal statistics and every pairwise
Kendall's $\tau$ -- direct edges and dependence induced through the graph
alike -- closely. The Gaussian-everywhere baseline, sampling the identical
graph with a single fixed correlation on every edge, instead produces a
uniform $\tau \approx 0.33$ regardless of an edge's true, and deliberately
varied, strength. On held-out data (library fitted on a training draw, compared with 100
independent test draws) the mean absolute error of the per-edge Kendall's
$\tau$ is 0.0084 for the learned models versus
0.159 for the baseline and
0.0114 for a fitted Gaussian copula; the
unbiased U-space energy distance to the true data is
-0.00007 (95\% CI
$[-0.00038, 0.00027]$, i.e.\ not significantly different from
zero) versus 0.0038, and a permutation test (run on the first
ten seeds) rejects equality of the generated and the true distribution in
0/10 seeds for the learned models versus 4/10 for the baseline. Because the
synthetic SCM is linear-Gaussian apart from Student-$t$ noise on one node,
the fitted Gaussian copula is nearly as good as the full family set here;
the advantage of richer families shows on Sachs (Section~\ref{sec:rq3}). This is the cleanest possible demonstration of what the
learned-dependency approach buys over a generic baseline: the graph
structure alone is not enough to make a sampler realistic, and a method that
ignores per-edge dependence strength gets every edge equally wrong unless
they all happen to share the same true strength by coincidence.

\section{RQ3: Does the choice of copula family set matter, and is that dataset-dependent?}
\label{sec:rq3}
Table~\ref{tab:sweep} reports the full hyperparameter sweep -- Gaussian-only
versus the full parametric family set, crossed with AIC versus BIC
selection -- for the two-parent conditional model on all five datasets.

\begin{table}[ht]
\caption{Conditional-vine AIC and mean $|\tau|$ error (in-sample) under Gaussian-only versus full parametric family selection, for each dataset's two-parent node. Asia's AIC- and BIC-selected fits are shown separately since they differ; every other dataset selects identical families under both criteria.}
\label{tab:sweep}
\centering
\small
\begin{tabular}{llrrl}
\toprule
\textbf{Dataset} & \textbf{Family set} & \textbf{AIC} & \textbf{Mean $|\tau|$ err.} & \textbf{Families selected} \\
\midrule
Sachs         & Gaussian-only & $-7{,}769.2$  & 0.0379 & gaussian, gaussian, gaussian \\
Sachs         & parametric    & $-13{,}649.8$ & 0.0490 & joe, joe, clayton \\
\addlinespace
UCI diabetes  & Gaussian-only & $-339.9$      & 0.0318 & gaussian, gaussian, gaussian \\
UCI diabetes  & parametric    & $-373.3$      & 0.0264 & frank, frank, clayton \\
\addlinespace
Synthetic SCM & Gaussian-only & $-6{,}267.7$  & 0.00272 & gaussian, gaussian, gaussian \\
Synthetic SCM & parametric    & $-6{,}265.2$  & 0.00264 & gaussian, clayton, gaussian \\
\addlinespace
Asia (AIC-sel.)& parametric   & $-2{,}224.8$  & 0.0127 & frank, clayton, clayton \\
Asia (BIC-sel.)& parametric   & $-2{,}219.5$  & 0.0128 & frank, indep, clayton \\
Asia          & Gaussian-only & $-1{,}563.3$  & 0.0328 & gaussian, gaussian, gaussian \\
\addlinespace
Insurance     & Gaussian-only & $-5{,}337.3$  & 0.0233 & gaussian, gaussian, gaussian \\
Insurance     & parametric    & $-5{,}849.4$  & 0.00617 & frank, frank, frank \\
\bottomrule
\end{tabular}
\end{table}

Allowing the full parametric family set, rather than restricting selection
to Gaussian only, improves the fitted conditional vine's AIC on Sachs from
$-7{,}769$ to $-13{,}650$ -- the largest improvement of any dataset -- consistent
with the fact that Sachs' pairwise fits are almost all non-Gaussian (Joe 26,
Clayton 14, Student-$t$ 11, Gumbel 2, Frank 1 and Gaussian 1 of the 55 pairs
on the pooled data). Insurance and, once the
sweep was extended to discrete data, Asia show a similarly large AIC
improvement (roughly $-512$ and $-660$ points respectively) and a
correspondingly large drop in mean $|\tau|$ error, and in both cases the
selected families are entirely non-Gaussian (Frank, Clayton). The picture
is different for the two most linear datasets. UCI diabetes shows a modest
AIC improvement ($-339.9 \to -373.3$) with a real, if smaller, drop in
$\tau$ error. The synthetic SCM, which was in fact generated with linear-Gaussian
relationships, shows the full family set buying essentially
nothing: AIC moves from $-6{,}267.7$ to $-6{,}265.2$, a change in the
\emph{wrong} direction. This is not a parameter-count penalty -- the
Clayton copula chosen for one edge has one parameter, like the Gaussian it
replaces. The reason is that the vine is selected tree by tree
\parencite{dissmann2013}: the second
tree is fitted to pseudo-observations computed from the first tree's
copulas, so a family that is locally better for one first-tree edge does not
guarantee a better joint fit. This is a useful negative result in its own
right: family-set flexibility is not a free lunch, and on data that
genuinely is close to linear-Gaussian, sequential (greedy) family selection
can very slightly worsen the joint fit.

Figure~\ref{fig:aic-improvement} visualizes the AIC improvement from
allowing the full parametric family set, expressed as $\Delta\text{AIC} =
\text{AIC}_{\text{parametric}} - \text{AIC}_{\text{Gaussian-only}}$ (more
negative bars indicate a larger improvement). Figure~\ref{fig:tau-error}
shows the corresponding change in mean $|\tau|$ error, which -- unlike raw
AIC -- is directly comparable across datasets since it does not depend on
sample size or dimensionality. The two figures agree for Asia and Insurance
(large AIC improvement, large drop in $\tau$ error) and for the synthetic
SCM (neither changes), but not for Sachs: the largest AIC improvement comes
with a \emph{higher} $\tau$ error (0.038 to 0.049). AIC scores the joint
density of the three variables, while the $\tau$ error looks only at the two
child--parent pairs, so AIC is not a reliable proxy for dependence
reconstruction and both are reported. The $\tau$ errors in this table are
in-sample and based on a single conditional target per dataset; differences
in the third decimal (e.g.\ 0.00272 versus 0.00264 on the synthetic SCM)
are within sampling noise and are not interpreted.

\begin{figure}[ht]
\centering
\begin{tikzpicture}
\begin{axis}[
  width=0.95\textwidth, height=5.5cm,
  ybar, bar width=14pt,
  symbolic x coords={Sachs,UCI diabetes,Synthetic SCM,Asia,Insurance},
  xtick=data, x tick label style={rotate=20, anchor=east, font=\small},
  ylabel={$\Delta$AIC (parametric $-$ Gaussian-only)},
  ymin=-6800, ymax=600,
  nodes near coords, every node near coord/.append style={font=\tiny},
  enlarge x limits=0.15,
  axis lines*=left,
]
\addplot[fill=gray!60] coordinates {
  (Sachs,-5880.6) (UCI diabetes,-33.4) (Synthetic SCM,2.5) (Asia,-661.5) (Insurance,-512.1)
};
\end{axis}
\end{tikzpicture}
\caption{AIC improvement from the full parametric family set relative to Gaussian-only, per dataset. A positive bar (synthetic SCM) means the richer family set actually made AIC slightly worse.}
\label{fig:aic-improvement}
\end{figure}

\begin{figure}[ht]
\centering
\begin{tikzpicture}
\begin{axis}[
  width=0.95\textwidth, height=5.7cm,
  ybar, bar width=10pt,
  symbolic x coords={Sachs,UCI diabetes,Synthetic SCM,Asia,Insurance},
  xtick=data, x tick label style={rotate=20, anchor=east, font=\small},
  ylabel={Mean $|\tau|$ error},
  ymin=0, ymax=0.045,
  enlarge x limits=0.15,
  legend style={at={(0.5,1.12)}, anchor=south, legend columns=2, font=\small},
  axis lines*=left,
]
\addplot[fill=gray!30] coordinates {
  (Sachs,0.0379) (UCI diabetes,0.0318) (Synthetic SCM,0.00272) (Asia,0.0328) (Insurance,0.0233)
};
\addplot[fill=gray!70] coordinates {
  (Sachs,0.0490) (UCI diabetes,0.0264) (Synthetic SCM,0.00264) (Asia,0.0127) (Insurance,0.00617)
};
\legend{Gaussian-only, Parametric}
\end{axis}
\end{tikzpicture}
\caption{Mean $|\tau|$ error under Gaussian-only versus full parametric family selection, per dataset. Sachs is a partial exception: the richer families fit the joint likelihood far better (Figure~\ref{fig:aic-improvement}) while showing a slightly higher $\tau$-recovery error, since AIC optimizes the vine's joint density rather than any single marginal-tau reconstruction target.}
\label{fig:tau-error}
\end{figure}

On the three continuous datasets and on Insurance, AIC and BIC select
identical models. Asia is the one dataset where they diverge: for the
parent--parent edge \emph{bronc}--\emph{either} of the first tree, AIC,
whose penalty per parameter (2) is lighter than BIC's ($\log n \approx 9.0$
for $n = 8{,}000$), keeps a one-parameter Clayton copula, while BIC prefers
the zero-parameter independence family (Table~\ref{tab:sweep}, ``frank,
clayton, clayton'' versus ``frank, indep, clayton''). The two selections are close enough in mean $\tau$ error
(0.0127 versus 0.0128) that either is a reasonable choice in practice, but
the fact that they disagree at all is informative: model-selection
ambiguity of this kind appeared only on a discrete dataset in this sweep,
and specifically on the smallest one, where a single edge's evidence for
one extra parameter is genuinely marginal rather than clear-cut.

\subsection{Indegree-Cap Sweep: Does Correctness Hold at Three to Six Parents?}
A separate sweep checks the sampler itself, independent of any specific
dependency library, as the indegree cap is raised from the core two-parent
target to six parents. Each run samples a node with $k$ independent root
parents using the Gaussian baseline with correlation $\rho_k$ between the
child and every parent and measures Kendall's $\tau$ across the $k$ edges;
each cap is repeated with ten seeds. For a Gaussian copula $\tau =
\frac{2}{\pi}\arcsin\rho$. With $\rho = 0.5$ for every $k$, as in the
original version of this experiment, the correlation matrix is positive
definite only for $k \le 3$: at $k = 4$ it is singular (the child is a
deterministic function of its parents) and for $k \ge 5$ it is not a valid
correlation matrix at all, so the original four-parent result tested a
degenerate model. The sampler now caps the correlation at $\rho_k =
\min(0.5, \sqrt{0.9/k})$, which leaves $k \le 3$ unchanged.
Table~\ref{tab:indegree} shows the observed mean $\tau$ within 0.006 of the
theoretical value at every cap, with no fallback draw.

\begin{table}[ht]
\caption{Indegree-cap sweep: baseline sampler with $\rho_k = \min(0.5, \sqrt{0.9/k})$, $n = 5{,}000$, ten seeds per cap. Mean $\tau$ $\pm$ standard deviation over seeds; min/max over all edges and seeds. For seed 0, caps 2 and 3 reproduce the original single-run values (mean $\tau$ 0.334 and 0.327).}
\label{tab:indegree}
\centering
\begin{tabular}{lcccccc}
\toprule
\textbf{Cap} & \textbf{$\rho_k$} & \textbf{$\tau$ theory} & \textbf{Mean $\tau$} & \textbf{Min} & \textbf{Max} & \textbf{Time (s)} \\
\midrule
2 & 0.500 & 0.333 & $0.339 \pm 0.006$ & 0.315 & 0.361 & $0.023 \pm 0.003$ \\
3 & 0.500 & 0.333 & $0.335 \pm 0.006$ & 0.316 & 0.356 & $0.039 \pm 0.004$ \\
4 & 0.474 & 0.315 & $0.314 \pm 0.004$ & 0.282 & 0.335 & $0.066 \pm 0.006$ \\
5 & 0.424 & 0.279 & $0.278 \pm 0.004$ & 0.257 & 0.298 & $0.081 \pm 0.012$ \\
6 & 0.387 & 0.253 & $0.253 \pm 0.003$ & 0.235 & 0.274 & $0.126 \pm 0.011$ \\
\bottomrule
\end{tabular}
\end{table}

Figure~\ref{fig:indegree-tau} plots the deviation of the observed from the
theoretical $\tau$ and Figure~\ref{fig:indegree-time} the wall-clock time,
which grows roughly linearly with the cap over this range (sampling only;
the baseline vine is refitted in every run). Neither shows any sign of
degrading correctness.

\begin{figure}[ht]
\centering
\begin{tikzpicture}
\begin{axis}[
  width=0.7\textwidth, height=5.5cm,
  xlabel={Indegree cap}, ylabel={$\bar\tau_{\text{observed}} - \tau_{\text{theory}}$},
  xtick={2,3,4,5,6}, xmin=1.7, xmax=6.3,
  ymin=-0.02, ymax=0.02,
  axis lines*=left,
]
\addplot[thick, dashed, gray] coordinates {(1.7,0) (6.3,0)};
\addplot[
  only marks, mark=*, mark size=2pt, black,
  error bars/.cd, y dir=both, y explicit
] coordinates {
  (2,0.0058) +- (0,0.0056)
  (3,0.0016) +- (0,0.0055)
  (4,-0.0008) +- (0,0.0041)
  (5,-0.0011) +- (0,0.0038)
  (6,-0.0006) +- (0,0.0031)
};
\end{axis}
\end{tikzpicture}
\caption{Deviation of the mean observed Kendall's $\tau$ from the theoretical Gaussian-copula value at each indegree cap (error bars: standard deviation over ten seeds).}
\label{fig:indegree-tau}
\end{figure}

\begin{figure}[ht]
\centering
\begin{tikzpicture}
\begin{axis}[
  width=0.7\textwidth, height=5.5cm,
  xlabel={Indegree cap}, ylabel={Elapsed time (s)},
  xtick={2,3,4,5,6}, xmin=1.7, xmax=6.3,
  ymin=0, ymax=0.15,
  axis lines*=left,
]
\addplot[thick, mark=square*, mark size=2pt,
  error bars/.cd, y dir=both, y explicit] coordinates {
  (2,0.0225) +- (0,0.0031) (3,0.0393) +- (0,0.0036) (4,0.0657) +- (0,0.0061)
  (5,0.0811) +- (0,0.0116) (6,0.1256) +- (0,0.0106)
};
\end{axis}
\end{tikzpicture}
\caption{Wall-clock time to fit the baseline vine and sample $n=5{,}000$ rows, per cap (mean $\pm$ standard deviation over ten seeds).}
\label{fig:indegree-time}
\end{figure}

The baseline check validates the sampling machinery. For \emph{learned}
vines, conditional vines were fitted for every multi-parent node of
Insurance (indegree 2--3) and Alarm (indegree 2--4) on a training draw and
used to sample each child given the parents of 10 independent test draws:
no fallback draw occurred at any indegree (0 in total), including the
40 child--parent edges of Alarm's four-parent nodes. The fidelity of the
reproduced child--parent dependence is governed by the discrete-data issue
of Section~\ref{sec:rq4} rather than by the number of parents: with
jittered fitting the mean absolute Cram\'er's V error of the child--parent
edges is 0.20/0.13 on Insurance (indegree 2/3) and
0.37/0.21/0.17 on Alarm (indegree 2/3/4); with native discrete fitting
it drops to 0.055/0.066 and
0.24/0.14/0.065 respectively.

Finally, the clipping guard of the conditional sampler was swept from
$\epsilon = 10^{-10}$ to $10^{-2}$ on the Sachs and Asia targets of
Table~\ref{tab:sweep}. Results are identical up to $\epsilon = 10^{-6}$
(the default); larger values start to bias the reproduced dependence (mean
$|\tau|$ error on Sachs 0.049 at $10^{-6}$, 0.050 at $10^{-3}$, 0.057 at
$10^{-2}$), and no setting triggered a fallback.

\section{RQ4: Does the pipeline extend to discrete and mixed data without redesign?}
\label{sec:rq4}
\label{sec:discrete-results}
Reconstructing the full Asia network end-to-end with the unmodified sampler
reproduces the true category frequencies closely: in the original
single-run evaluation the mean total-variation distance (TVD) over nodes is
$0.0033$ for the Gaussian baseline and $0.0024$ for the learned model, both
without any fallback draw. On Insurance the baseline reached only $0.0574$
against $0.0045$ for the learned model. The gap is concentrated at the
multi-parent nodes (mean TVD of the baseline 0.004 at roots, 0.010 at
one-parent, 0.062 at two-parent and 0.091 at three-parent nodes): the
baseline's multi-parent model is a Gaussian vine fitted with mutually
\emph{independent} parents, but in the generated graph the parents are
dependent, so the child's generated uniforms are no longer uniform and its
categories are mis-weighted. This is the compatibility issue described in
Section~\ref{sec:sampler}, not an accumulation of error along the graph;
with margin calibration the baseline's TVD drops to the level of the
learned model (Table~\ref{tab:rq4}). In this sense -- can the continuous
pipeline be pointed at discrete data and produce something sane -- the
extension succeeds without changing the fitting or sampling code.

Per-edge dependence is where jittering shows its cost. In the original
evaluation, the learned model's mean absolute error in Cram\'er's V over the
Asia edges is even slightly \emph{worse} than the naive baseline's ($0.2662$
versus $0.2439$; normalized mutual information $0.2172$ versus $0.2188$);
Insurance shows the same pattern ($0.1859$ versus $0.1815$, NMI $0.1360$
versus $0.1324$). The worst-matched Asia edges are \emph{lung}--\emph{either}
(true Cram\'er's V $0.898$, simulated $0.300$), \emph{either}--\emph{xray}
($0.753 \to 0.307$) and \emph{tub}--\emph{either} ($0.431 \to 0.112$); on
Insurance, \emph{MakeModel}--\emph{Antilock} ($0.638 \to 0.213$),
\emph{Age}--\emph{SeniorTrain} ($0.728 \to 0.307$) and
\emph{VehicleYear}--\emph{Airbag} ($0.841 \to 0.432$). The mechanism is
explained below, and its remedy, native discrete fitting, is evaluated on
held-out data in Table~\ref{tab:rq4}.

Figure~\ref{fig:asia-network} shows the full Asia network structure, so
the specific edges named throughout this section can be located within the
graph rather than treated as an unstructured list of variable-name pairs.
The edges examined in Table~\ref{tab:jitter} below span the network's
skew spectrum deliberately: \emph{smoke}$\to$\emph{bronc} sits near the
graph's root with a roughly balanced split, while \emph{lung}$\to$\emph{either}
sits one step before the highly skewed \emph{either} node, whose own
definition as a logical OR of two upstream binary variables makes it the
most extreme case of the skew this section quantifies.

\begin{figure}[ht]
\centering
\begin{tikzpicture}[
  node/.style={draw, circle, minimum size=8mm, font=\scriptsize},
  arr/.style={-{Latex[length=1.8mm]}, thick}
]
\node[node] (asia) at (0,1.2) {asia};
\node[node] (smoke) at (0,-1.2) {smoke};
\node[node] (tub) at (2,1.2) {tub};
\node[node] (lung) at (2,-0.4) {lung};
\node[node] (bronc) at (2,-2) {bronc};
\node[node] (either) at (4,0.4) {either};
\node[node] (xray) at (6,1.2) {xray};
\node[node] (dysp) at (6,-1.2) {dysp};

\draw[arr] (asia) -- (tub);
\draw[arr] (smoke) -- (lung);
\draw[arr] (smoke) -- (bronc);
\draw[arr] (tub) -- (either);
\draw[arr] (lung) -- (either);
\draw[arr] (either) -- (xray);
\draw[arr] (either) -- (dysp);
\draw[arr] (bronc) -- (dysp);
\end{tikzpicture}
\caption{The Asia network \parencite{lauritzen1988}: 8 binary variables, indegree at most 2. \emph{either} is a deterministic logical OR of \emph{tub} and \emph{lung}, making it the most heavily skewed node in the graph.}
\label{fig:asia-network}
\end{figure}

The jittering transform dilutes measured concordance relative to the raw
data's Kendall's $\tau_b$, and the dilution is severe for skewed or rare
binary variables. The mechanism is an identity rather than an
approximation \parencite{denuit2005,genest2007}. After independent
jittering, a pair of observations that is tied in $x$ (or $y$) in the raw
data is ordered at random and contributes zero expected concordance; all
other pairs keep their order. Hence $\tau$ of the jittered data equals, in
expectation, Kendall's $\tau_a = (n_c - n_d)/\binom{n}{2}$ of the raw data,
and because $\tau_b = (n_c - n_d)/\sqrt{(n_0 - n_x)(n_0 - n_y)}$ with $n_0 =
\binom{n}{2}$ and $n_x, n_y$ the numbers of pairs tied in $x$ and $y$,
\[
\mathbb{E}\bigl[\tau_{\text{jittered}}\bigr] = \tau_a = \tau_b(\text{raw}) \cdot \sqrt{(1-t_x)(1-t_y)},
\qquad t_x = n_x/n_0,\; t_y = n_y/n_0 .
\]
Table~\ref{tab:jitter} shows this on real Asia edges, including ordinary
ones and not only the near-deterministic \emph{either = tub OR lung} node
that first drew attention to the effect; the remaining differences are the
noise of a single jitter draw. A copula fitted to jittered data therefore
\emph{correctly} fits a diluted dependence -- the input signal genuinely is
weaker than the raw categorical dependence -- which is why a naive
Gaussian baseline sometimes lands, by chance, closer to the raw value on an
individual edge than the jitter-based model does. The remedy is to avoid
jittering altogether and fit the copula with the native discrete likelihood
(Section~\ref{sec:discrete-method}); Table~\ref{tab:rq4} and the paragraph
after it quantify the improvement.

\begin{table}[ht]
\caption{Predicted versus observed jittered Kendall's $\tau$ on four Asia edges spanning a range of category skew, verified against the dilution formula.}
\label{tab:jitter}
\centering
\begin{tabular}{lcccc}
\toprule
\textbf{Edge} & \textbf{Raw $\tau_b$} & \textbf{Jittered $\tau$ (observed)} & \textbf{Predicted} & \textbf{Abs.\ error} \\
\midrule
smoke $\to$ lung   & 0.196 & 0.037 & 0.044 & 0.007 \\
smoke $\to$ bronc  & 0.307 & 0.153 & 0.153 & 0.001 \\
either $\to$ dysp  & 0.195 & 0.041 & 0.048 & 0.006 \\
lung $\to$ either  & 0.898 & 0.101 & 0.099 & 0.002 \\
\bottomrule
\end{tabular}
\end{table}

These four edges were chosen deliberately to span a range of category
skew rather than to show the effect at its most extreme: \emph{smoke
$\to$ bronc} joins two roughly balanced variables (49/51 and 55/45) and
shows only mild dilution; in \emph{either $\to$ dysp} the parent is skewed
93.5/6.5 (\emph{dysp} is 56.5/43.5) and the dilution is strong; \emph{lung
$\to$ either} is skewed on both sides (94.7/5.3 and 93.5/6.5) and shows
near-total dilution despite a very high raw $\tau_b$ of 0.898. A fifth
edge, \emph{asia $\to$ tub}, was fitted but is excluded from the table: both
variables are about 99\% ``no'' (98.8/1.2 and 98.7/1.3),
so the raw dependence signal is small enough to be dominated by
BIF-resampling noise rather than by the dilution mechanism itself, and
including it would muddy the comparison rather than sharpen it.

Figure~\ref{fig:dilution-scatter} plots predicted against observed jittered
$\tau$ for the four edges on a $y=x$ reference line; as expected from the
identity, every point lies on the line up to the noise of one jitter draw.

\begin{table}[ht]
\caption{RQ4 on held-out data: the library is fitted once on a training draw and each of 50 independent test draws is compared with the generated data (mean over seeds, 95\% bootstrap CI in brackets). TVD: mean category-frequency total-variation distance over discrete nodes; CV: mean absolute Cram\'er's V error over discrete--discrete edges. ``cal.'' = margin calibration.}
\label{tab:rq4}
\centering
\small
\begin{tabular}{llcc}
\toprule
\textbf{Network} & \textbf{Generator} & \textbf{TVD} & \textbf{CV error} \\
\midrule
Asia & fixed-$\rho$ baseline & 0.0040 & 0.240 [0.238, 0.242] \\
 & fixed-$\rho$ baseline, cal. & 0.0034 & 0.240 [0.238, 0.242] \\
 & learned, jittered (thesis method) & 0.0043 & 0.266 [0.264, 0.267] \\
 & learned, native discrete & 0.0058 & 0.140 [0.139, 0.142] \\
 & learned, native discrete, cal. & 0.0034 & 0.139 [0.138, 0.141] \\
 & CPT refit (true model class) & 0.0047 & 0.015 [0.014, 0.016] \\
\addlinespace
Insurance & fixed-$\rho$ baseline & 0.0579 & 0.182 [0.181, 0.182] \\
 & fixed-$\rho$ baseline, cal. & 0.0062 & 0.176 [0.176, 0.177] \\
 & learned, jittered (thesis method) & 0.0079 & 0.187 [0.187, 0.188] \\
 & learned, native discrete & 0.0107 & 0.061 [0.060, 0.062] \\
 & learned, native discrete, cal. & 0.0062 & 0.058 [0.057, 0.059] \\
 & CPT refit (true model class) & 0.0073 & 0.011 [0.011, 0.012] \\
\addlinespace
Alarm & fixed-$\rho$ baseline & 0.0095 & 0.310 [0.309, 0.311] \\
 & fixed-$\rho$ baseline, cal. & 0.0055 & 0.314 [0.313, 0.314] \\
 & learned, jittered (thesis method) & 0.0063 & 0.345 [0.344, 0.345] \\
 & learned, native discrete & 0.0088 & 0.187 [0.186, 0.187] \\
 & learned, native discrete, cal. & 0.0055 & 0.188 [0.187, 0.189] \\
 & CPT refit (true model class) & 0.0065 & 0.012 [0.012, 0.013] \\
\addlinespace
Mixed SCM & fixed-$\rho$ baseline & 0.0055 & 0.250 [0.247, 0.253] \\
 & fixed-$\rho$ baseline, cal. & 0.0029 & 0.249 [0.246, 0.252] \\
 & learned, jittered (thesis method) & 0.0060 & 0.240 [0.237, 0.244] \\
 & learned, native discrete & 0.0059 & 0.008 [0.007, 0.010] \\
 & learned, native discrete, cal. & 0.0029 & 0.008 [0.006, 0.010] \\
\bottomrule
\end{tabular}
\end{table}

Two experiments quantify the native-discrete remedy. First, for every edge
of Asia and Insurance a single bivariate copula was fitted on a training
draw, once to jittered pseudo-observations and once with the native
discrete likelihood, and pairs simulated from it were mapped back through
the discrete quantile function and compared with an independent test draw.
The mean absolute Cram\'er's V error falls from 0.263 to 0.023 on the eight
Asia edges (Wilcoxon $p = 0.016$) and from 0.177 to 0.045 on the 52
Insurance edges ($p < 10^{-4}$); on \emph{lung}--\emph{either}, for
example, the jittered copula reproduces a Cram\'er's V of 0.30 and the
native one 0.89 (test data 0.92). Second, Table~\ref{tab:rq4} reconstructs
whole networks. Native fitting roughly halves the edge-dependence error on
Asia (0.266 to 0.140) and Alarm
(0.345 to 0.187), cuts it by two thirds on Insurance
(0.187 to 0.061) and almost removes it on the mixed
SCM (0.240 to 0.0084); every one of these differences, and
the difference to the fixed-$\rho$ baseline, has a paired Wilcoxon
$p < 10^{-4}$ over the 50 seeds. On the three discrete networks uncalibrated native
fitting slightly increases the category-frequency error (TVD), which
margin calibration removes. On the purely discrete networks, refitting the conditional
probability tables, which is the true model class of the data, remains
clearly better (Cram\'er's V error between 0.011 and 0.015); the
copula generator's advantage is that the same machinery handles continuous,
discrete and mixed data and allows dependencies to be transplanted between
graphs, which CPTs do not.

\begin{figure}[ht]
\centering
\begin{tikzpicture}
\begin{axis}[
  width=0.65\textwidth, height=6.5cm,
  xlabel={Predicted $\tau_{\text{jittered}}$}, ylabel={Observed $\tau_{\text{jittered}}$},
  xmin=0, xmax=0.18, ymin=0, ymax=0.18,
  axis lines*=left,
  legend pos=north west, legend style={font=\small},
]
\addplot[thick, dashed, gray, domain=0:0.18] {x};
\addlegendentry{$y=x$}
\addplot[only marks, mark=*, mark size=2.3pt, black] coordinates {
  (0.044,0.037) (0.153,0.153) (0.048,0.041) (0.099,0.101)
};
\addlegendentry{Asia edges}
\end{axis}
\end{tikzpicture}
\caption{Predicted versus observed jittered Kendall's $\tau$ for the four Asia edges in Table~\ref{tab:jitter}. All four points lie within 0.007 of the $y=x$ line.}
\label{fig:dilution-scatter}
\end{figure}

\subsection{Mixed Continuous/Discrete Data}
Table~\ref{tab:mixed} reports every edge of the dedicated mixed synthetic
SCM (continuous columns $x_1, x_3, x_4$; discrete columns $x_2, x_5$; edges
$x_1{\to}x_2$, $x_1{\to}x_4$, $x_2{\to}x_5$, $x_3{\to}x_4$, $x_4{\to}x_5$).
Both marginal families are reproduced closely on their own: the three
continuous columns show a Kolmogorov--Smirnov statistic no larger than
$0.013$, and the two discrete columns show a total-variation distance no
larger than $0.0064$.

\begin{table}[ht]
\caption{Per-edge Kendall's $\tau$ (true vs.\ simulated, ordinal-numeric view) on the mixed synthetic SCM, by edge type.}
\label{tab:mixed}
\centering
\begin{tabular}{llccc}
\toprule
\textbf{Edge} & \textbf{Type} & \textbf{$\tau_{\text{true}}$} & \textbf{$\tau_{\text{sim}}$} & \textbf{Abs.\ error} \\
\midrule
$x_1 \to x_4$ & continuous--continuous & 0.460 & 0.470 & 0.010 \\
$x_3 \to x_4$ & continuous--continuous & 0.363 & 0.348 & 0.015 \\
$x_1 \to x_2$ & continuous--discrete   & 0.643 & 0.547 & 0.096 \\
$x_4 \to x_5$ & continuous--discrete   & 0.504 & 0.308 & 0.197 \\
$x_2 \to x_5$ & discrete--discrete     & 0.644 & 0.404 & 0.241 \\
\bottomrule
\end{tabular}
\end{table}

The ordering in Table~\ref{tab:mixed} (jittered fitting) is not a
coincidence: dependence error grows from continuous--continuous ($\le
0.015$), to one-side-discrete ($0.096$--$0.197$), to fully
discrete--discrete ($0.241$, corroborated by a Cram\'er's V drop from
$0.683$ true to $0.428$ simulated on the same pair). This is what the
identity above predicts: an edge with one continuous and one discrete
endpoint picks up a single $\sqrt{1-t}$ factor from the discrete side,
while an edge between two discrete endpoints picks up the full
$\sqrt{(1-t_x)(1-t_y)}$ product. With native discrete fitting the
discrete--discrete error disappears (Table~\ref{tab:rq4}: mean Cram\'er's V
error 0.007). The
mixed dataset therefore does more than confirm that Phase 1 and Phase 2
results combine as expected -- it isolates edge type as the variable that
predicts how much dilution to expect, holding the rest of the graph fixed.

Figure~\ref{fig:mixed-edges} shows the same five edges as a bar chart
grouped by type, which makes the step pattern easier to see at a glance
than the table alone.

\begin{figure}[ht]
\centering
\begin{tikzpicture}
\begin{axis}[
  width=0.85\textwidth, height=5.5cm,
  ybar, bar width=16pt,
  symbolic x coords={$x_1\to x_4$,$x_3\to x_4$,$x_1\to x_2$,$x_4\to x_5$,$x_2\to x_5$},
  xtick=data, x tick label style={font=\small},
  ylabel={Abs.\ error in $\tau$},
  ymin=0, ymax=0.27,
  nodes near coords, every node near coord/.append style={font=\tiny},
  enlarge x limits=0.12,
  axis lines*=left,
]
\addplot[fill=gray!25] coordinates {
  ($x_1\to x_4$,0.010) ($x_3\to x_4$,0.015) ($x_1\to x_2$,0.096) ($x_4\to x_5$,0.197) ($x_2\to x_5$,0.241)
};
\end{axis}
\end{tikzpicture}
\caption{Dependence error by edge, grouped visually by type: the two leftmost bars are continuous--continuous, the middle two are one-side-discrete, and the rightmost is discrete--discrete -- a monotonically increasing step that matches the dilution formula's prediction.}
\label{fig:mixed-edges}
\end{figure}


\subsection{Conditional Independences of the Supplied Graph}
\label{sec:rq4-ci}
Reproducing marginals and edge dependence does not guarantee that generated
data obey the conditional independences the supplied graph implies, and for
discrete variables this does not follow from the construction either. On
the copula scale a child depends on its parent only through the parent's
uniform $u_B$, so $A \perp C \mid u_B$ holds in a chain or fork by
construction. A discrete $B$, however, reveals only the category slice in
which $u_B$ lies, while $A$ and $C$ both still depend on where $u_B$ lies
inside that slice; given the observed category of $B$, they can therefore
remain dependent. This is the situation analysed by \textcite{sun2024}: when
the middle variable is observed only in discretized form, a conditional
independence that holds for the underlying continuous variables need not
hold for the observations, and standard tests on the observations then
reject it. Whether the generated data preserve the graph's independences
therefore has to be checked. The check below deliberately tests them on the
observed scale, because that is the scale on which discovery algorithms
use the data; the networks' own conditional probability tables satisfy
them there exactly.

\emph{Protocol.} Every induced chain ($A \to B \to C$), fork
($A \leftarrow B \to C$) and collider ($A \to B \leftarrow C$) of the true
DAGs of Asia, Insurance, Alarm and the mixed SCM, with $A$ and $C$
non-adjacent, was taken as a supplied three-node graph (at most ten per
type and network, chosen at random; 77 in total). For each, the generators
learned their building blocks from the training draw (seed 10{,}000) of the
three columns and generated $n$ rows ($8{,}000$; $5{,}000$ for the mixed SCM)
per evaluation seed. On every sample the implied independence was tested --
$A \perp C \mid B$ for chains and forks, $A \perp C$ for colliders -- with a
likelihood-ratio ($G$) test on the contingency table (continuous $A$ or $C$
binned into five quantile classes, which keeps the null hypothesis valid)
or, for a continuous $B$, the generalised covariance measure
\parencite{shah2020}. The implied dependence ($A$--$C$ for chains and forks,
$A$--$C \mid B$ for colliders) was tested in the same way. Under a true
independence about 5\% of the tests reject. Two references calibrate the
tests: refitted CPTs on the three-node graph (discrete triples), whose data
satisfy the implied independences exactly, and real data wherever the full
network implies the same statement ($d$-separation).

\begin{table}[ht]
\caption{Share of tests (\%) that reject the independence implied by the supplied three-node graph ($A \perp C \mid B$ for chains and forks, $A \perp C$ for colliders), $\alpha = 0.05$; about 5\% are expected when it holds. Number of evaluation seeds after the network, number of triples after the structure. CPT: refitted conditional probability tables; Cop.: the learned copula generator with jittered (J) or native discrete (N) fitting. --: not applicable (no CPT refit for continuous variables; real data only where the full network implies the statement).}
\label{tab:ci-checks}
\centering
\small
\begin{tabular}{llcccccc}
\toprule
\textbf{Network} & \textbf{Structure} & \textbf{Real} & \textbf{CPT} & \textbf{Baseline} & \textbf{Cop.\ (J)} & \textbf{Cop.\ (N)} & \textbf{Forest} \\
\midrule
Asia (20) & chain (7) & 4 & 6 & 80 & 17 & 44 & 3 \\
 & fork (2) & 12 & 5 & 95 & 15 & 70 & 5 \\
 & collider (2) & 5 & 2 & 8 & 8 & 8 & 8 \\
\addlinespace
Insurance (10) & chain (10) & -- & 7 & 56 & 18 & 29 & 7 \\
 & fork (10) & 0 & 9 & 40 & 18 & 44 & 9 \\
 & collider (10) & 10 & 9 & 9 & 9 & 9 & 9 \\
\addlinespace
Alarm (10) & chain (10) & 2 & 6 & 80 & 42 & 74 & 7 \\
 & fork (10) & 5 & 5 & 97 & 32 & 78 & 14 \\
 & collider (10) & 6 & 7 & 8 & 8 & 8 & 8 \\
\addlinespace
Mixed SCM (20) & chain (3) & -- & -- & 22 & 22 & 40 & 7 \\
 & fork (1) & 0 & -- & 5 & 10 & 15 & 5 \\
 & collider (2) & 5 & -- & 5 & 5 & 0 & 5 \\
\bottomrule
\end{tabular}
\end{table}

Table~\ref{tab:ci-checks} shows that the copula generators do \emph{not}
reliably preserve the independences of chains and forks whose middle
variable is discrete. Pooled over these 50 triples, the implied
$A \perp C \mid B$ is rejected in 25\% of the tests for the learned
generator with jittering, 56\% with native discrete fitting and 72\%
for the fixed-$\rho$ baseline, against 7\% for refitted CPTs, 5\%
for real data and 7\% for random-forest conditionals; the rejection
rate exceeds one half for 18\%, 52\% and 68\% of these triples,
respectively, and for none with CPTs or forests. The violations are small
in size (median conditional mutual information below 0.003 nats) but
systematic. Native fitting, which reproduces edge dependence far better
(Section~\ref{sec:rq4}), violates them more often than jittering, which is
consistent with the mechanism above: the stronger the latent dependence on
$u_B$, the more information about $A$ and $C$ remains within a category of
$B$. The forest does not have this problem, because it draws a
discrete child from the parent's category alone. With a continuous middle
variable, all generators stay close to the nominal level (8\% for the
jittered copula, 8\% for the forest), and the marginal independence of
the two parents of a collider holds for every generator (8\% of the tests
reject for each, against 6\% for real data).

The implied dependences are present in real data and CPT-refit data in
94\% and 90\% of the tests on the discrete networks, but in only
50\% for the jittered copula generator, whose dependence is diluted
(Section~\ref{sec:rq4}), against 84\% with native fitting and 90\%
for the forest. On discrete data, the copula generator therefore errs in
both directions that matter to a constraint-based algorithm: it weakens
dependences the graph implies and adds conditional dependences the graph
excludes. This is consistent with the spurious and missing edges PC finds
on its data for Insurance and Alarm (Section~\ref{sec:closing-loop}).

\section{RQ5: Are the learned dependencies genuinely reusable across datasets and graphs?}
\label{sec:composition}
Every result so far reconstructs a dataset's \emph{own} graph. RQ5 tests
reuse in structurally different contexts, in two experiments
(\path{gcda/experiments/composition.py}).

\subsection{A graph built from two datasets}
The Sachs and Asia libraries were merged
(\texttt{DependencyLibrary.merge}) and the graph \emph{praf} $\to$
\emph{pmek}, \emph{praf} $\to$ \emph{dysp} was sampled, which exists in
neither dataset: \emph{praf} and \emph{pmek} keep their Sachs marginals and
the learned \emph{praf}--\emph{pmek} copula, \emph{dysp} keeps its Asia
marginal, and the cross-dataset edge \emph{praf} $\to$ \emph{dysp} reuses
the copula Asia learned for \emph{bronc} $\to$ \emph{dysp}, \emph{dysp}'s
strongest parent (a Frank copula with $\tau = 0.339$; the graph
specification format is given in Appendix~\ref{sec:cli}). The run uses strict
mode, i.e.\ no baseline model is involved. Over 30 seeds of $n = 10{,}000$
rows: the marginals are reproduced (KS 0.008 for \emph{praf} and
\emph{pmek}, TVD 0.005 for \emph{dysp}); the \emph{praf}--\emph{pmek}
dependence matches its copula ($\tau = 0.669$ versus 0.666); and on the
copula scale the transplanted edge reproduces its source copula to within
sampling error ($\tau = 0.338 \pm 0.005$ versus 0.339). On the observable scale the
binary child attenuates it to $\tau_a = 0.243$ (versus $\tau_a = 0.334$ for
\emph{bronc}--\emph{dysp} in Asia, where \emph{both} variables are
discrete), which illustrates that a transplanted dependence is defined at
the copula level and its observable strength depends on the marginals it
is attached to. In the earlier demonstration of this capability the
cross-dataset edge had silently used the fixed $\rho = 0.5$ baseline
because no binding mechanism existed; the binding (and the CLI support for
it) were added as part of this work.

\subsection{Novel graphs over the Sachs variables}
200 random DAGs over the eleven Sachs variables (maximum indegree 2, none
of them the consensus graph) were sampled with four generators: learned
pairwise copulas plus a conditional vine fitted for every two-parent node
of the graph (``fitted vine''); learned pairwise copulas with two-parent
nodes composed from the edges' pairwise copulas (``composed'',
Section~\ref{sec:sampler}); the same with Gaussian-only families (``fitted
Gaussian''); and the fixed-$\rho$ baseline. Table~\ref{tab:novel} reports,
for $n = 5{,}000$, how well PC recovers the \emph{specified} graph, how
faithfully the one-parent edges reproduce their library copula, marginal
fidelity and varsortability \parencite{reisach2021}.

\begin{table}[ht]
\caption{RQ5: 200 random DAGs over the Sachs variables (mean over DAGs). SHD and skeleton $F_1$ of PC (Fisher-$z$, $\alpha=0.05$) against the specified DAG; $\Delta\tau$: mean $|\tau_{\text{generated}} - \tau_{\text{copula}}|$ over one-parent edges; KS: mean marginal KS distance to the real data; VS: varsortability (fraction of edges along which the marginal variance increases; a linear-Gaussian SCM with random weights on the same DAGs gives 0.92). Copulas fitted on the pooled (P) or observational (O) Sachs rows.}
\label{tab:novel}
\centering
\small
\begin{tabular}{llccccc}
\toprule
\textbf{Rows} & \textbf{Generator} & \textbf{SHD} & \textbf{$F_1$} & \textbf{$\Delta\tau$} & \textbf{KS} & \textbf{VS} \\
\midrule
P & fitted vine      & 11.73 & 0.77 & 0.007 & 0.013 & 0.50 \\
P & composed         & 11.66 & 0.78 & 0.007 & 0.012 & 0.49 \\
P & fitted Gaussian  & 11.37 & 0.81 & 0.007 & 0.012 & 0.49 \\
P & fixed $\rho$     & 11.85 & 0.79 & 0.008 & 0.017 & 0.52 \\
\addlinespace
O & fitted vine      & 13.47 & 0.33 & 0.007 & 0.012 & 0.49 \\
O & composed         & 13.46 & 0.34 & 0.007 & 0.012 & 0.49 \\
O & fitted Gaussian  & 13.60 & 0.30 & 0.007 & 0.012 & 0.49 \\
O & fixed $\rho$     & 11.47 & 0.79 & 0.008 & 0.017 & 0.51 \\
\bottomrule
\end{tabular}
\end{table}

Three findings stand out. First, the building blocks transfer to graphs
that never existed: one-parent edges reproduce the $\tau$ of the copula they
were built from to within 0.007 on average (95\% CI up to 0.008), and the
pairwise-composed two-parent model shows no significant difference in SHD
from a vine fitted specifically for that parent set (paired difference
$+0.07$, $p = 0.75$, on the pooled and
$+0.01$, $p = 0.90$, on the observational copulas). As in
Section~\ref{sec:closing-loop}, this bounds the difference rather than
showing that the two are equivalent.
Second, \emph{how hard} the generated graph is to recover is inherited from
the real data: with the observational copulas, most of which are weak, PC
recovers only a third of the skeleton, whereas the fixed-$\rho$ baseline,
which makes every edge equally strong, is easy (SHD
2.00 lower than with the fitted vines, 95\% CI
$[1.49, 2.50]$, $p < 10^{-4}$). With the pooled copulas, which are
stronger, the generators differ by less than half an SHD unit (fitted
vine versus baseline $p = 0.33$; versus the fitted Gaussian copula,
which is slightly easier to recover, $+0.36$, Wilcoxon
$p = 0.048$, $t$-test $p = 0.11$). Generated benchmarks therefore reflect
the dependence strengths of the source data rather than an arbitrary
simulation choice -- which is the point of the method, but it also means
that ``usable'' must be read as ``faithful to the specified graph and the
learned dependencies'', not ``easy to recover''. Third, the varsortability
of all copula-generated data is about 0.5, i.e.\ marginal variances carry
no information about the causal order, in contrast to
0.92 for a linear-Gaussian SCM on the same graphs;
benchmarks generated this way cannot be gamed by sorting variables by
variance. The generated marginals deviate slightly from the real ones
(mean KS about 0.013); with margin calibration this falls to 0.004
without materially changing any other metric.

\section{RQ6: Is the generated data actually useful for causal discovery?}
\label{sec:closing-loop}
Every check up to this point validates statistical realism. RQ6 asks
whether a structure-learning algorithm can use the generated data as it
can use real data. Table~\ref{tab:shd} gives the original four-seed result;
Table~\ref{tab:shd-heldout} the held-out evaluation of
Section~\ref{sec:reproduce-seeds}.

\begin{table}[ht]
\caption{Original protocol: mean (standard deviation) SHD between the PC algorithm's recovered graph and the true DAG, over 4 seeds (the library was fitted on the seed-1 draw of the synthetic SCM, which is also one of the evaluation seeds).}
\label{tab:shd}
\centering
\begin{tabular}{lcc}
\toprule
\textbf{Source} & \textbf{Continuous (synthetic SCM)} & \textbf{Discrete (Asia)} \\
\midrule
True data & 1.25 (0.50) & 6.50 (0.58) \\
Baseline (Gaussian everywhere) & 1.00 (0.00) & 8.25 (1.26) \\
Learned (T4 models) & 1.00 (0.00) & 6.75 (1.26) \\
\bottomrule
\end{tabular}
\end{table}

With four seeds, the Asia difference between the learned generator and
the baseline (1.5 SHD) was not statistically significant (paired $t$-test
$p = 0.22$). Table~\ref{tab:shd-heldout} repeats the experiment on held-out
data with 100 seeds, stronger baselines and two larger networks, and
Table~\ref{tab:shd-ges} repeats it with GES.

\begin{table}[ht]
\caption{Held-out RQ6 with PC: mean SHD to the true DAG [95\% bootstrap CI] and mean number of edges PC finds (true number of edges in parentheses), 100 seeds per graph.}
\label{tab:shd-heldout}
\centering
\small
\begin{tabular}{llcc}
\toprule
\textbf{Graph} & \textbf{Source} & \textbf{SHD} & \textbf{Edges found} \\
\midrule
Synth.\ SCM (4) & true data & 1.13 [1.06, 1.21] & 4.1 \\
 & fixed-$\rho$ baseline & 1.13 [1.06, 1.21] & 4.1 \\
 & fitted Gaussian copula & 1.14 [1.04, 1.28] & 4.1 \\
 & learned (jittered) & 1.15 [1.04, 1.29] & 4.1 \\
\addlinespace
Asia (8) & true data & 6.37 [6.27, 6.47] & 5.7 \\
 & fixed-$\rho$ baseline & 8.18 [7.87, 8.51] & 12.5 \\
 & fitted Gaussian copula & 7.68 [7.53, 7.84] & 5.9 \\
 & learned (jittered) & 6.56 [6.38, 6.76] & 7.7 \\
 & learned (native discrete) & 8.58 [8.43, 8.74] & 8.0 \\
 & CPT refit & 6.21 [6.13, 6.29] & 5.3 \\
\addlinespace
Insurance (52) & true data & 23.19 [22.64, 23.74] & 41.3 \\
 & fixed-$\rho$ baseline & 46.56 [45.70, 47.46] & 65.8 \\
 & fitted Gaussian copula & 41.34 [40.64, 42.06] & 33.7 \\
 & learned (jittered) & 37.39 [36.52, 38.24] & 36.2 \\
 & learned (native discrete) & 60.32 [59.69, 60.93] & 51.9 \\
 & CPT refit & 22.39 [21.95, 22.85] & 41.6 \\
\addlinespace
Alarm (46) & true data & 9.17 [8.75, 9.64] & 43.9 \\
 & fixed-$\rho$ baseline & 31.55 [30.86, 32.26] & 59.6 \\
 & fitted Gaussian copula & 37.50 [36.87, 38.13] & 36.0 \\
 & learned (jittered) & 30.58 [29.74, 31.41] & 43.6 \\
 & learned (native discrete) & 52.37 [51.68, 53.07] & 63.9 \\
 & CPT refit & 9.14 [8.75, 9.56] & 44.0 \\
\bottomrule
\end{tabular}
\end{table}

\begin{table}[ht]
\caption{Held-out RQ6 with GES (BIC score for the SCM, BDeu for the discrete networks): mean SHD to the true DAG [95\% bootstrap CI] and mean number of edges found, 100 seeds per graph.}
\label{tab:shd-ges}
\centering
\small
\begin{tabular}{llcc}
\toprule
\textbf{Graph} & \textbf{Source} & \textbf{SHD} & \textbf{Edges found} \\
\midrule
Synth.\ SCM (4) & true data & 1.05 [1.00, 1.14] & 4.0 \\
 & fixed-$\rho$ baseline & 1.04 [1.00, 1.12] & 4.0 \\
 & fitted Gaussian copula & 1.07 [1.01, 1.17] & 4.0 \\
 & learned (jittered) & 1.07 [1.01, 1.17] & 4.0 \\
\addlinespace
Asia (8) & true data & 3.10 [3.01, 3.23] & 7.4 \\
 & fixed-$\rho$ baseline & 10.25 [9.93, 10.56] & 11.2 \\
 & fitted Gaussian copula & 8.07 [8.00, 8.14] & 4.7 \\
 & learned (jittered) & 6.37 [5.99, 6.76] & 6.5 \\
 & learned (native discrete) & 7.96 [7.88, 8.04] & 7.3 \\
 & CPT refit & 3.10 [3.03, 3.20] & 7.2 \\
\addlinespace
Insurance (52) & true data & 22.59 [21.81, 23.40] & 43.2 \\
 & fixed-$\rho$ baseline & 36.14 [35.63, 36.66] & 57.1 \\
 & fitted Gaussian copula & 41.41 [40.72, 42.11] & 27.8 \\
 & learned (jittered) & 33.92 [33.64, 34.20] & 30.7 \\
 & learned (native discrete) & 49.58 [48.95, 50.22] & 47.4 \\
 & CPT refit & 20.46 [20.04, 20.88] & 43.6 \\
\addlinespace
Alarm (46) & true data & 11.29 [10.08, 12.55] & 45.5 \\
 & fixed-$\rho$ baseline & 24.60 [24.13, 25.08] & 52.0 \\
 & fitted Gaussian copula & 42.73 [42.15, 43.25] & 27.0 \\
 & learned (jittered) & 36.32 [35.26, 37.32] & 32.9 \\
 & learned (native discrete) & 38.59 [37.90, 39.29] & 59.7 \\
 & CPT refit & 9.98 [8.77, 11.23] & 45.7 \\
\bottomrule
\end{tabular}
\end{table}

On the continuous SCM all sources sit at the floor of about one error
(the one edge whose orientation is not identifiable from observational
data) under both algorithms, and no two sources differ significantly.

\emph{Asia, PC.} The learned generator's data supports PC significantly
better than the fixed-$\rho$ baseline (SHD difference
$-1.62$, 95\% CI $[-1.98, -1.27]$, Wilcoxon
$p < 10^{-4}$) and than a fitted Gaussian copula (difference
$-1.12$, 95\% CI $[-1.34, -0.89]$, $p < 10^{-4}$), so the advantage is due to the learned
non-Gaussian dependence, not merely to learned strengths. Its difference
from real data is not significant (difference
$+0.19$, 95\% CI $[-0.01, 0.41]$, Wilcoxon $p = 0.15$). This does not show
that the two are equivalent; it shows that any excess error of the
generated data is at most about 0.4 edges, the upper confidence bound. Refitted CPTs are
slightly better still ($+0.35$, $p = 0.003$). The baseline's
extra error comes from spurious edges: PC finds 12.5 edges on its data
against 7.7 on the learned generator's and 5.7 on real data (the true
graph has 8).

\emph{Asia, GES.} With GES the ranking of the generators is the same --
learned (6.37) before the fitted Gaussian copula
(8.07) and the fixed-$\rho$ baseline (10.25), $p < 10^{-4}$ for both
comparisons -- but GES recovers the graph far better from real data
(3.10) than from the learned generator's data, so the small difference
from real data seen with PC does not carry over to a score-based method.

\emph{Insurance and Alarm.} On the two larger networks the learned
generator falls far short of real data under both algorithms (PC: SHD
37.4 versus 23.2 on Insurance and 30.6 versus 9.2 on Alarm; GES:
33.9 versus 22.6 and 36.3 versus 11.3; all $p < 10^{-4}$), while refitted
CPTs come within about two edges of real data. Relative to the Gaussian alternatives the picture is
mixed: with PC the learned generator beats both on Insurance
($-9.2$ and $-4.0$ SHD, $p < 10^{-4}$) and the fitted
Gaussian copula on Alarm ($-6.9$), but is not significantly different from the
fixed-$\rho$ baseline on Alarm ($-0.97$, Wilcoxon
$p = 0.21$); with GES it beats both on Insurance but is \emph{worse} than the
fixed-$\rho$ baseline on Alarm ($+11.7$ SHD, $p < 10^{-4}$), where it
yields too sparse graphs (32.9 edges found against 45.5 on real data).

\emph{Native discrete fitting}, which reproduces the edge dependence far
better (Section~\ref{sec:rq4}), makes both algorithms \emph{worse} on every
discrete network (PC on Asia 8.58, Insurance 60.3, Alarm
52.4); on Alarm it also produces many more adjacencies than exist
(63.9 edges found with PC against 46 true edges), whereas on Asia and
Insurance the number of edges found is about right but many of them are
wrong. Section~\ref{sec:rq4-ci} shows one reason directly: on the observed
scale, native fitting violates the conditional independences of chains and
forks with a discrete middle variable in 56\% of the tests, against 25\%
with jittering and 7\% for refitted CPTs, so the tests PC and GES rely on
find dependences the graph excludes. A further plausible contribution is
that native fitting also reproduces near-deterministic relations such as
Asia's \emph{either} = \emph{tub} OR \emph{lung}, which violate the
faithfulness assumption the algorithms rely on. Statistical fidelity and usefulness for a given
discovery algorithm are therefore related but not interchangeable, and
producing data that is as usable for discovery as real data on large
discrete networks remains an open problem for copula conditionals.

\subsection{Random-Forest Conditionals}
\label{sec:rq6-forest}
Table~\ref{tab:rf} repeats the held-out RQ6 evaluation with one further
source: every non-root node sampled from a random-forest conditional
(Section~\ref{sec:forest-method}) fitted on the same training draw as the
copula library. The other sources are regenerated with the same code and
seeds and are identical, row by row, to those of
Tables~\ref{tab:shd-heldout} and~\ref{tab:shd-ges}.

\begin{table}[ht]
\caption{Held-out RQ6 with random-forest conditionals: mean SHD over 100 seeds, and the share of the gap between the fixed-$\rho$ baseline and real data that the copula generator (learned, jittered) and the forest close. Above $100\%$, discovery is easier than on real data.}
\label{tab:rf}
\centering
\small
\begin{tabular}{llccccccc}
\toprule
 & & \multicolumn{5}{c}{\textbf{SHD}} & \multicolumn{2}{c}{\textbf{Gap closed}} \\
\cmidrule(lr){3-7} \cmidrule(lr){8-9}
\textbf{Graph} & \textbf{Alg.} & \textbf{Baseline} & \textbf{Copula} & \textbf{Forest} & \textbf{CPT refit} & \textbf{Real} & \textbf{Copula} & \textbf{Forest} \\
\midrule
Asia & PC & 8.18 & 6.56 & 5.24 & 6.21 & 6.37 & $90\%$ & $162\%$ \\
Insurance & PC & 46.56 & 37.39 & 22.90 & 22.39 & 23.19 & $39\%$ & $101\%$ \\
Alarm & PC & 31.55 & 30.58 & 9.91 & 9.14 & 9.17 & $4\%$ & $97\%$ \\
\addlinespace
Asia & GES & 10.25 & 6.37 & 3.01 & 3.10 & 3.10 & $54\%$ & $101\%$ \\
Insurance & GES & 36.14 & 33.92 & 19.91 & 20.46 & 22.59 & $16\%$ & $120\%$ \\
Alarm & GES & 24.60 & 36.32 & 8.79 & 9.98 & 11.29 & $-88\%$ & $119\%$ \\
\bottomrule
\end{tabular}
\end{table}

On Insurance and Alarm, which the copula generator could not handle, the
forest brings both algorithms close to real data. With PC its difference
from real data is not significant on Insurance (difference
$-0.29$, 95\% CI $[-1.04, 0.48]$, Wilcoxon
$p = 0.63$; i.e.\ within about one edge either way) and small but significant on Alarm
($+0.74$, 95\% CI $[0.15, 1.33]$). With GES it overshoots: its
data yields a \emph{lower} SHD than real data on Insurance
($-2.68$, $p < 10^{-4}$) and Alarm
($-2.50$, $p = 0.0089$), which is a mismatch as
well, if a smaller one than the copula generator's. On Asia the forest
shows no significant difference from real data with GES ($-0.09$,
95\% CI $[-0.22, 0.01]$, $p = 0.57$) but
overshoots with PC ($-1.13$, 95\% CI $[-1.37, -0.88]$), where the copula
generator was closer to real data. On the continuous SCM, whose edges
are all monotone, the forest does not differ significantly from the copula
generator. Against the copula generator the forest lowers the SHD on every
discrete network under both algorithms (PC: $-1.32$ on Asia,
$-14.49$ on Insurance, $-20.67$ on Alarm; all
$p < 10^{-4}$).

The edge dependence explains the difference. The mean error of Cram\'er's
$V$ over the true edges falls from 0.268, 0.185 and 0.344 for the
copula generator to 0.019, 0.016 and 0.014 for the forest (Asia, Insurance,
Alarm), close to the CPT refit (0.016, 0.014, 0.013).
A model without the copula's monotonicity and simplifying assumptions
reproduces both the non-monotone edges and the interactions, and with them
the conditional independences the discovery algorithms test. On the
discrete networks the forest is, however, no better than refitting the
CPTs, the networks' own model class; its advantage over a CPT refit is
that the same code applies to continuous and mixed variables. Its
disadvantage relative to the copula library is the one noted in
Section~\ref{sec:forest-method}: it does not provide transferable building
blocks.

\chapter{Discussion}
\label{chap:discussion}

\section{What the Bugs Reveal About the Method, Not Just the Code}
\label{sec:bugs-discussion}
Several of the bugs found during development are worth revisiting
collectively rather than only as isolated fixes, because they share a
pattern: each one was invisible to generic testing and visible only to a
check informed by what the numbers \emph{should} mean statistically. None
of them would have been caught by a unit test that only checks output
shapes or types.

\subsection{The Pseudo-Observation Scrambling Bug}
Section~\ref{sec:bug-case-study-1} covers this bug's mechanics in full;
the point worth adding here is \emph{why} it evaded detection for as long
as it did. The function's signature (Listing~\ref{lst:bug1-before}) was
not wrong in a way any type checker or shape assertion could flag -- it
returned an array of the correct length, correct dtype, and correct range
$(0,1)$, just in the wrong row order. Every individual pairwise fit that
consumed this output ran successfully and produced a plausible-looking
single number. What made the bug visible was not inspecting any one
fit's output, but computing many fits at once and noticing that their
\emph{distribution} was wrong: real variable pairs from a real dataset
essentially never all show $\tau \approx 0.999$, which is what happens
when two independently sorted (and therefore functionally identical)
columns are compared to each other. This is the clearest illustration in
the whole project of aggregate, cross-fit sanity checking catching
something no single-fit test could.

\subsection{The Independence-Family Reshape Bug}
Chapter~\ref{chap:implementation}'s second case study is a different
failure mode: not wrong output, but a crash, and a crash that was entirely
deterministic given the right input -- yet absent for the whole of Phase 1.
The lesson here is about test \emph{coverage} of input space rather than
output correctness: independence was selected for a few Phase-1 pairs (six
in UCI diabetes, one in the synthetic SCM), but none of them was a
one-parent edge of a graph that was sampled, so the \texttt{indep} branch
of the deserialization code was simply never
exercised, by any test or any demo, during the whole of Phase 1. Asia -- with several
genuinely independent binary pairs -- was the first dataset to hit it. The
regression test added afterward does not just fix this one case; it forces
the independence family specifically into the test matrix, which is the
only way to guarantee this class of bug (works for every family present in
the data so far, breaks on a family that happens not to appear yet) cannot
recur unnoticed for some future dataset either.

\subsection{The Sweep Validation-Guard Bug}
A third case, mentioned only briefly in earlier chapters, fits the same
pattern from a slightly different angle. \path{dependence_strength_sweep}'s
correlation override was written to apply to a node's single incoming
pairwise copula, but the first demo call happened to target a node that
actually had two parents in that DAG, so the override silently had nothing
to act on and Kendall's $\tau$ stayed constant across every value tested
in the sweep. A flat line across a sweep that is supposed to show a
monotonic trend is exactly the same kind of implausible-summary-statistic
signal as the $\tau \approx 0.999$ case above, just inverted: too little
variation rather than too much. The fix was not merely to correct the
target node, but to add an explicit validation guard that raises clearly
whenever the target child does not have exactly one parent, converting a
silent no-op into a loud, immediate failure the next time the same mistake
is made.

\subsection{A Common Thread}
The standard-normal baseline bug -- found independently in both the
continuous evaluation script and, later, its discrete counterpart --
extends this pattern once more: a baseline comparison that was supposed to
isolate dependence-model quality instead used a generic marginal instead
of each dataset's real one, conflating two different sources of error into
a single number. What all four bugs share is that they hide behind a
result that looks superficially reasonable -- a single fitted parameter,
a single sweep row, a single baseline comparison -- and only become
visible when checked against either (a) a broader ensemble of similar
results that should vary, or (b) an expected qualitative behavior (a
monotonic trend, a clean isolation between two quantities) that a
one-off inspection would not think to verify. The practical implication
for anyone extending this pipeline is that its test suite needs
statistical sanity checks with a known expected direction or magnitude,
not just structural checks, and that an implausibly clean or implausibly
flat result is itself worth investigating before being reported as a
finding.

\section{The Audit Pass}
\label{sec:audit}
Midway through Phase 2, the implemented scope was checked directly against
the original task plan rather than assumed complete, and this surfaced five
concrete gaps: the ordinal-category-order default described in
Section~\ref{sec:discrete-method}; a mixed continuous/discrete pipeline that
had only ever been unit-tested, never run end-to-end; a hyperparameter
sweep that never actually covered discrete data (fixing this is what first
revealed AIC and BIC diverging, Section~\ref{sec:rq3}); a missing baseline
comparison for the discrete pipeline, whose addition is what led directly to
discovering the jitter-dilution effect in Section~\ref{sec:rq4}; and a
three-parent closed-form verification that had been done ad hoc during
development but never committed as a regression test. All five were
concrete, fixable gaps rather than open design questions, which suggests the
audit was worth doing as a discrete step rather than trusting that
task-by-task completion implies scope-level completion.

A second, external-style audit of this document against the code (the
record is part of the repository, \path{thesis_audit/}) found that several
descriptions had drifted from the implementation -- most importantly the
vine structure and Rosenblatt order (Sections~\ref{sec:conditional}
and~\ref{sec:sampler}), a cross-dataset demonstration whose transplanted
edge silently used the baseline, and a parent-order-sensitive model lookup
that made the same graph produce different data depending on the order in
which its edges were listed. It also found that a four-parent baseline was
degenerate and that the downstream result rested on four seeds, one of
which overlapped the training data. All of these were fixed in code, tested,
and re-evaluated in Chapter~\ref{chap:results}.

\section{Reproducibility Across Environments}
Running the pipeline on a cluster with a different Python version and
freshly resolved (rather than pinned) dependency versions reproduced every
continuous result exactly, since those only depend on \texttt{numpy}'s
seeded generator, which is bit-stable across platforms for the operations
used here. The discrete sweep results did not reproduce, and tracking down
why surfaced a genuine bug rather than an environment quirk: the sweep
script's call into the categorical continuization step omitted a fixed
seed, so it drew fresh operating-system entropy for jitter randomness on
every run, while every other call site in the codebase passed a fixed seed
and the sweep's own continuous rows never exercised that code path at all.
This is a useful cautionary example independent of this specific project:
a bug that never manifests within a single environment can still be a real
bug, and cross-environment reruns are a legitimate way to find it -- not a
compliance exercise, but a genuine second layer of testing that single-machine
development cannot provide by itself. After the fix, a third environment
(Python 3.13 with newer releases of \texttt{pyvinecopulib},
\texttt{pgmpy}, \texttt{causal-learn} and \texttt{numpy}) reproduced the
sweep table bit-for-bit (maximum absolute AIC difference $10^{-11}$) and the
four-seed SHD table exactly. That environment is pinned in
\path{requirements-lock.txt}, every output carries a provenance file, and a
continuous-integration job re-runs the sweep and compares it with the
committed table on every commit. On the Lichtenberg cluster the
dependency libraries were then rebuilt from scratch and every original
evaluation script re-run: the sweep table again agreed to $10^{-11}$ in AIC
with identical families, and every evaluation reproduced its committed
values, with a single exception at the fourth decimal (the Insurance
baseline's mean Cram\'er's V error, $0.1816$ instead of $0.1815$), which
the original code shows as well when run with the current library
versions.

\section{Limitations}
Jittering-based continuization measures Kendall's $\tau_a$ instead of
$\tau_b$ and therefore under-reproduces dependence for skewed categorical
variables (Section~\ref{sec:rq4}). The native discrete likelihood removes
most of this loss, but copulas of discrete data are not unique
\parencite{genest2007}, so what is transplanted from a discrete dataset is
a latent dependence that is only identified up to that ambiguity; on a
purely discrete network, refitting the conditional probability tables
remains the better generator for that network's own graph.

Reusing a multi-parent model in a new graph is exact only if the parents'
joint distribution matches the one the model was fitted under; otherwise
the child's marginal drifts. The sampler detects this and margin
calibration repairs the marginal, but calibration preserves only the ranks
of the generated child, not its conditional distribution. The pairwise
composition of multi-parent nodes assumes conditional independence of the
parents given the child in copula space.

The downstream evaluation uses PC and GES on four known graphs with 100
seeds each. It supports the specific claims of
Section~\ref{sec:closing-loop}, not general statements about other
algorithm families, larger graphs, or supervised structure-learning
methods trained on generated data. Generated data is only as close to
reality as its building blocks: with GES on Asia, and with both algorithms
on Insurance and Alarm, the graph is recovered markedly less well from
the copula generator's data than from real data, so for discrete networks
of this size the copula generator should not be used as a stand-in for
real data. Where a generator is reported as not significantly different
from real data, this bounds the difference by the confidence interval but
does not establish equivalence; no equivalence margin was specified in
advance. The generated discrete data also do not reliably preserve the conditional independences of the supplied graph: with a discrete middle variable, chain and fork independences are violated by the copula generators (Section~\ref{sec:rq4-ci}), so on the observed scale their data are not Markov with respect to the supplied graph. Random-forest conditionals
close most of the discovery gap (Section~\ref{sec:rq6-forest}), but are fitted per
node, parent set and dataset and therefore not reusable building blocks;
they were evaluated only on the known graphs of RQ6.

\section{Threats to Validity}
The Sachs file pools nine experimental conditions; the RQ1 results are
therefore reported on the pooled rows and on the 853-row observational
condition, whose identity was verified row by row against an independent
copy of the observational data. RQ1 compares a ranking with a consensus network
that is itself an estimate. The Sachs measurements are rounded (91\% ties
in \emph{praf}); the fitted copulas were insensitive to the tie-handling
rule, but empirical and copula-implied $\tau$ differ (0.616 versus 0.666
for \emph{praf}--\emph{pmek}).

Each held-out experiment fits its library on a single training draw, so
the variance due to library fitting is not resampled. The discrete networks
are sampled from bnlearn CPTs via pgmpy, and their deterministic or
near-deterministic nodes (Asia's \emph{either}) violate faithfulness, which
affects every source including real data. Finally, the composition study
uses random DAGs over one dataset's variables; transplanting dependencies
between semantically unrelated variables is a modelling choice of the user
whose plausibility the tool cannot judge.

The conditional-independence check of Section~\ref{sec:rq4-ci} uses at most
ten randomly chosen triples per structure and network and, for Insurance and
Alarm, only ten seeds. Its $G$-test is slightly liberal on sparse
contingency tables: on Insurance, data from refitted CPTs, for which the
tested independences hold exactly, are rejected in 7--9\% of the tests
rather than 5\%. The check is deliberately run on the observed variables,
the scale on which discovery algorithms see the data, not on the latent
continuous variables a discretization-aware test
\parencite{sun2024} would target.

\chapter{Conclusion and Future Work}
\label{chap:conclusion}

\section{Summary}
This thesis set out to build a tool that generates realistic synthetic data
for an arbitrary user-specified graph by learning local dependencies from
real data rather than assuming a fixed generic dependence model, and to
verify on held-out data that the result is both statistically faithful and
usable for causal discovery. The strongest fitted dependencies coincide
with known causal edges (RQ1); sampling reproduces the true, non-uniform
per-edge dependence where a generic baseline cannot (RQ2); flexible family
selection matters in a dataset-dependent way, and the sampler is correct up
to six parents with its baseline and four with learned vines (RQ3); the
pipeline extends to discrete and mixed data, where native discrete fitting
removes most of the dependence loss that jittering causes (RQ4); learned
building blocks transfer to graphs that never existed, including across
datasets, and yield benchmarks that are not variance-sortable (RQ5); and on
Asia, PC recovers the true graph significantly better from the learned
generator's data than from a fixed-correlation baseline or a fitted
Gaussian copula, with no significant difference from real data, whereas GES on Asia, and both
algorithms on the larger Insurance and Alarm networks, fall clearly short of
real data (RQ6). The downstream claim is therefore limited to the small
Asia network with PC: on Insurance and Alarm, data from the copula generator
are not a substitute for real data when discovery algorithms are evaluated
(PC: SHD 37.4 versus 23.2 and 30.6 versus 9.2; GES: 33.9 versus 22.6 and
36.3 versus 11.3). On discrete data, the copula generators also violate the conditional independences of chains and forks with a discrete middle variable (rejected in 25\%--56\% of the tests, against 7\% for refitted CPTs and for random-forest conditionals; Section~\ref{sec:rq4-ci}).
Random-forest conditionals fitted on the same data close most of the
discovery gap on the known graphs, which is consistent with the explanation
that it stems from non-monotone and interaction dependence that copula
conditionals cannot represent; they were not evaluated on novel graphs.

\section{Future Work}
The main open problem is the gap to real data on larger discrete networks
(Section~\ref{sec:closing-loop}). Random-forest conditionals close most of
it on the known graphs (Section~\ref{sec:rq6-forest}) but are not reusable
building blocks. A hybrid that keeps copulas for monotone edges and for
dependence transferred to new graphs, and uses nonparametric conditionals
only where an edge is non-monotone or parents interact, would combine both
properties; it would also have to be evaluated on novel graphs, which the
prototype was not. Beyond that,
three directions follow from the results. First, generated benchmarks should be used for what motivated
them: training and evaluating supervised or amortized structure-learning
methods on many generated (graph, dataset) pairs, and testing whether
performance on generated data predicts performance on real data. Second,
the compatibility of transplanted multi-parent models deserves a principled
treatment -- for instance re-estimating a vine's parent-dependence trees
from the parents actually generated, while keeping the child's conditional
copulas. Third, a systematic study of which cross-dataset compositions
produce plausible variables would turn a working capability into a
principled one.

\appendix
\chapter{Software Artifact}

\section{Repository Layout}
The codebase is organized as an installable package (\texttt{gcda/}) plus a
set of narrated, single-purpose scripts (\texttt{scripts/}) that mirror the
same operations as the composable \texttt{gcda} command-line interface,
described below. The repository is available at
\url{https://github.com/Ramya-ravi26/Graph-Conditioned-Causal-Data-Augmentation}.
Every result file in \texttt{outputs/} is accompanied by a
\texttt{.meta.json} file with the package versions, the arguments and the
commit of the development repository in which the result was produced.

\begin{table}[ht]
\caption{Top-level repository structure.}
\label{tab:repo-layout}
\centering
\small
\begin{tabular}{lp{9.5cm}}
\toprule
\textbf{Path} & \textbf{Contents} \\
\midrule
\texttt{gcda/data/} & Dataset loaders (incl.\ the bundled Sachs CSV as package data), dataset registry, preprocessing (continuous, jittered and native-discrete inputs). \\
\texttt{gcda/dependency/} & Pairwise and conditional copula fitting, pairwise-composed multi-parent models, random-forest conditionals, \texttt{DependencyLibrary} (merge, validation). \\
\texttt{gcda/generation/} & \texttt{GraphSampler} (fallback counters, drift diagnostics, margin calibration) and the graph-specification layer used by the CLI. \\
\texttt{gcda/evaluation/} & Realism (incl.\ unbiased energy distance/MMD and permutation tests), graph-conditioning, discrete metrics, PC/GES + SHD. \\
\texttt{gcda/experiments/} & Hyperparameter sweep and the held-out, multi-seed experiments of Chapter~\ref{chap:results}. \\
\texttt{gcda/provenance.py} & Run metadata (package versions, git commit, arguments) written next to every output. \\
\texttt{gcda/cli.py} & The \texttt{gcda} command (build-library / generate / evaluate / sweep). \\
\texttt{scripts/} & Narrated entry points for each pipeline stage, \texttt{run\_experiments.py}, result merging and reproduction checks. \\
\texttt{configs/sweep.yaml} & Versioned definition of what the sweep varies. \\
\texttt{outputs/} & Versioned result artifacts the tables and figures are produced from. \\
\texttt{slurm/}, \texttt{docs/HPC\_RUNBOOK.md} & Job scripts and instructions for the large-seed runs on the Lichtenberg cluster. \\
\texttt{tests/} & Pytest suite mirroring the package layout, including every regression test referenced in Chapters~\ref{chap:method} and~\ref{chap:discussion}. \\
\texttt{thesis/}, \texttt{thesis\_audit/} & This document's source and the record of its corrections. \\
\texttt{LAB\_NOTEBOOK.md} & Chronological record of decisions, validations, and bugs found and fixed. \\
\bottomrule
\end{tabular}
\end{table}

\section{Command-Line Interface}
\label{sec:cli}
\begin{verbatim}
gcda build-library --dataset sachs --output lib_sachs.json
gcda build-library --dataset asia --discrete-mode native \
    --output lib_asia.json
gcda generate --dag graph.json --library lib_sachs.json \
    --library lib_asia.json --n 5000 --output data.csv [--strict]
gcda evaluate --true-csv true.csv --sim-csv data.csv --distances
gcda sweep [--wandb]
\end{verbatim}
\texttt{graph.json} accepts any DAG, not necessarily a source dataset's own
structure. In its simplest form it is an edge list,
\path{{"edges": [["a", "b"], ["a", "c"]]}}, whose nodes are looked up in
the dataset given by \texttt{--dataset}. Each node can instead be mapped
to a (dataset, variable) of any of the merged libraries, and each edge can
be bound to any learned pair copula (Section~\ref{sec:composition} uses
exactly this to build a graph from Sachs and Asia building blocks).
Multi-parent nodes use a fitted vine when one exists for that parent set
and are otherwise composed from the edges' pairwise copulas
(Section~\ref{sec:sampler}). Every modelling decision and every baseline
fallback is written to a JSON report next to the output, and
\texttt{--strict} turns any fallback into an error.

\section{Sweep Configuration}
Listing~\ref{lst:sweep-yaml} shows \texttt{configs/sweep.yaml}, the
versioned file the sweep (\path{gcda/experiments/sweep.py}) reads rather
than encoding these choices as constants -- the reason the sweep results in
Table~\ref{tab:sweep} and Table~\ref{tab:indegree} are reproducible from a
single committed file. Sample sizes and seeds of the datasets are fixed in
\path{gcda/experiments/sweep.py} (synthetic SCM $n=5{,}000$, Asia and
Insurance $n=8{,}000$); W\&B logging is offline unless
\texttt{WANDB\_MODE=online} is set.

\begin{lstlisting}[caption={configs/sweep.yaml (verbatim, comments removed).}, label={lst:sweep-yaml}]
family_sets:
  gaussian_only: [gaussian]
  parametric: [indep, gaussian, student, clayton, gumbel, frank, joe]
selection_criteria: [aic, bic]
targets:
  sachs:         {child: pmek, parents: [praf, PKA]}
  uci_diabetes:  {child: disease_progression, parents: [s5, bmi]}
  synthetic_scm: {child: x4, parents: [x1, x2]}
  asia:          {child: dysp, parents: [bronc, either]}
  insurance:     {child: DrivQuality, parents: [RiskAversion, DrivingSkill]}
indegree_sweep:
  caps: [2, 3, 4, 5, 6]
  baseline_rho: 0.5
  n: 5000
  seed: 0
  repeats: 10
clip_sweep:
  datasets: [sachs, asia]
  eps_values: [1.0e-10, 1.0e-6, 1.0e-3, 1.0e-2]
\end{lstlisting}

\section{Environment}
\label{sec:environment}
\texttt{pyproject.toml} declares the direct dependencies with lower bounds
(the oldest versions the results were produced with), the Python floor of
3.10 (needed for the \texttt{X | Y} union type hints used throughout the
codebase) and a single upper bound, \texttt{pgmpy<1.3}, because pgmpy 1.3
removes the bundled-network loader the discrete experiments use. In
addition, \texttt{requirements-lock.txt} pins the exact environment
(Python 3.13) in which every table of Chapter~\ref{chap:results} was
reproduced bit-for-bit, and every output file is accompanied by a
\texttt{.meta.json} file recording package versions and the git commit.
Table~\ref{tab:deps} lists the direct dependencies.

\begin{table}[ht]
\caption{Direct dependencies (lower bounds in \texttt{pyproject.toml}; exact versions in \texttt{requirements-lock.txt}).}
\label{tab:deps}
\centering
\begin{tabular}{ll}
\toprule
\textbf{Package} & \textbf{Role} \\
\midrule
\texttt{pyvinecopulib} & pairwise and vine copula fitting/sampling (Section~\ref{sec:pairwise-math}) \\
\texttt{networkx} & DAG representation and topological sort (Section~\ref{sec:sampler}) \\
\texttt{pgmpy} & bundled Asia/Insurance CPT sampling (Section~\ref{sec:discrete-method}) \\
\texttt{causal-learn} & PC algorithm for the downstream discovery check (Section~\ref{sec:closing-loop}) \\
\texttt{numpy}, \texttt{scipy} & numerical core \\
\texttt{pandas} & tabular data handling \\
\texttt{scikit-learn} & UCI diabetes loader, general utilities \\
\texttt{wandb} & sweep logging (offline by default) \\
\texttt{PyYAML} & sweep configuration \\
\texttt{matplotlib} & diagnostic plots \\
\texttt{pytest} & test suite \\
\bottomrule
\end{tabular}
\end{table}

\section{Reproducing These Results}
\label{sec:reproduce}
The results in Chapter~\ref{chap:results} are regenerated with the commands
below, run from the repository root after installing the locked environment
(Section~\ref{sec:environment}). Each command is followed by the results it
produces.

\medskip\noindent\textbf{Environment and tests.}
{\footnotesize
\begin{verbatim}
pip install -r requirements-lock.txt
pip install -e . --no-deps
pytest -q
\end{verbatim}
}

\noindent\textbf{Original pipeline} (RQ1--RQ6 as first reported;
Tables~\ref{tab:sweep}, \ref{tab:indegree}, \ref{tab:jitter}, \ref{tab:mixed}
and~\ref{tab:shd}).
{\footnotesize
\begin{verbatim}
python scripts/build_dependency_library.py
python scripts/generate_graph_data.py
python scripts/sweep_hyperparameters.py --output-dir /tmp/sweep
python scripts/check_reproduction.py /tmp/sweep/sweep_dependency_models.csv \
    outputs/sweep_dependency_models.csv
python scripts/build_discrete_library.py
python scripts/evaluate_discrete_generation.py
python scripts/evaluate_mixed_generation.py
python scripts/demo.py
python scripts/evaluate_causal_discovery.py
\end{verbatim}
}

\noindent\textbf{Held-out, multi-seed experiments}
(\path{scripts/run_experiments.py}; results in \path{outputs/experiments/}).
{\footnotesize
\begin{verbatim}
python scripts/run_experiments.py rq1
python scripts/run_experiments.py rq2 --seeds 100
python scripts/run_experiments.py rq4 --seeds 50 \
    --graphs asia insurance alarm mixed_scm
python scripts/run_experiments.py rq5
python scripts/run_experiments.py composition --n-dags 200 --seeds 30
python scripts/run_experiments.py rq6 --algorithm pc --seeds 100 \
    --graphs synthetic_scm asia insurance alarm
python scripts/run_experiments.py rq6 --algorithm ges --seeds 100 \
    --graphs synthetic_scm asia insurance alarm
python scripts/run_experiments.py discrete-native
python scripts/run_experiments.py learned-indegree --seeds 10
python scripts/run_experiments.py rf --rf-algorithms pc ges --seeds 100 \
    --graphs synthetic_scm asia insurance alarm
python scripts/run_experiments.py ci-checks --graphs asia mixed_scm --seeds 20
python scripts/run_experiments.py ci-checks --graphs insurance alarm --seeds 10
\end{verbatim}
}

\noindent The 100-seed and 200-graph runs were split into Slurm array jobs
on the Lichtenberg cluster (\path{slurm/*.slurm}, with
\path{scripts/merge_results.py} combining the tasks; see
\path{docs/HPC_RUNBOOK.md}). Their merged results are in
\path{outputs/experiments/hpc/}.

Every command above writes its numeric output to \texttt{outputs/} together with
a provenance file, so that a later step -- or a reader checking a specific
number against this thesis -- can consume a previous step's exact output
rather than re-deriving it by eye from terminal text. These artifacts are
versioned in the repository (\texttt{scripts/check\_reproduction.py}
compares a fresh sweep with them; the continuous-integration workflow runs
this check on every commit).

\section{Test Suite Reference}
Table~\ref{tab:test-suite} lists the regression tests referenced throughout
this thesis, organized by the module they cover and the specific bug or
property each one guards against, following the testing philosophy laid
out in Chapter~\ref{chap:implementation}: every test here asserts a
statistical property or an input-space edge case that was previously
missed, not just a structural shape or type check.

\begin{table}[ht]
\caption{Regression tests referenced in this thesis, by module.}
\label{tab:test-suite}
\centering
\begin{tabular}{lp{8.7cm}}
\toprule
\textbf{Test module} & \textbf{What it guards against} \\
\midrule
\texttt{test\_preprocessing.py} & Pseudo-observation round-trip on a permuted column (Section~\ref{sec:bug-case-study-1}); the jitter-dilution factor on a simulated skewed binary pair (Section~\ref{sec:rq4}). \\
\texttt{t3\_validation.py} (script) & Family/parameter/$\tau$ recovery on toy Gaussian, Student-$t$ and Clayton data with known ground truth (Section~\ref{sec:pairwise-math}); exits non-zero on failure. \\
\texttt{test\_pairwise.py} & Gaussian family/$\tau$ recovery; independence detection for independent columns. \\
\texttt{test\_conditional.py} & Two-parent conditional sampler against closed-form Gaussian theory (Section~\ref{sec:conditional}); the three-parent extension of the same closed-form check, added during the audit pass (Section~\ref{sec:audit}). \\
\texttt{test\_graph\_sampler.py} & Fallback on failing and on non-finite edge models; the independence-family $(0,0)$-shape reshape guard (Chapter~\ref{chap:implementation}); the $h$-function direction check by binning (Section~\ref{sec:conditional}). \\
\texttt{test\_library.py} & JSON save/load round trip of marginals, pairwise and conditional models. \\
\texttt{test\_graph\_conditioning.py} & \texttt{dependence\_strength\_sweep}'s one-parent validation guard (Section~\ref{sec:bugs-discussion}); edge removal and monotonicity checks. \\
\texttt{test\_remediation.py} & Regression tests for the corrections made after the audit: parent-order invariance of conditional lookups (library and CLI), library merge and namespace collisions, copula binding with orientation, cross-dataset CLI generation, strict mode, pairwise-composed conditionals against closed-form Gaussian theory, margin calibration and drift warnings, the valid baseline at four and five parents, per-node fallback counters, unbiased distances and permutation tests, library validation, marginal compression, parallel fitting, and the $\tau_a$ identity behind the dilution formula (Section~\ref{sec:rq4}). \\
\texttt{test\_forest.py} & Random-forest conditionals (Section~\ref{sec:forest-method}): a logical OR child and a non-monotone continuous child reproduced, integration with the graph sampler for one and two parents, no replay of training rows. \\
\texttt{test\_ci\_checks.py} & Chain/fork/collider detection and d-separation (Section~\ref{sec:rq4-ci}); the G-test and the generalised covariance measure hold their 5\% level under a true conditional independence, including step-shaped conditional means and repeated values of the conditioning variable. \\
\texttt{test\_datasets.py} & Loaders, ordinal category orders, the Sachs condition boundaries and observational subset, the consensus DAG (Mek has three parents). \\
\bottomrule
\end{tabular}
\end{table}

\printbibliography

\end{document}
